%----------------------------------------------
%	PACKAGES AND OTHER DOCUMENT CONFIGURATIONS
%----------------------------------------------
\documentclass[11pt, a4paper, oneside]{Thesis} % Paper size, default font size and one-sided paper

\graphicspath{{./Pictures/}} % Specifies the directory where pictures are stored

%\usepackage{pdfpages} % for including pdf files
\usepackage{float} % force figures and tables to be placed where they are referenced by using the paramter H.
\usepackage{color}
\usepackage{listings}
\usepackage{enumerate}

\usepackage{lscape} % for landscape oriented pages
\usepackage[square, numbers, comma, sort&compress]{natbib}
% Use the natbib reference package - read up on this to edit the reference style;
%if you want text (e.g. Smith et al., 2012) for the in-text references (instead of numbers), remove 'numbers' 

\definecolor{maroon}    {rgb}{.5, 0, 0}
\definecolor{green}     {rgb}{.1,.8,.1}
\definecolor{darkgreen} {rgb}{ 0,.5, 0}
\definecolor{lightgray} {rgb}{.9,.9,.9}
\definecolor{darkgray}  {rgb}{.4,.4,.4}

\definecolor{punct}     {rgb}{0.5,0,0}
\definecolor{delim}     {rgb}{0.08,0.41,0.69}
\definecolor{numb}      {rgb}{0.8,0.2,0.2}

\lstdefinelanguage{XML}
{
  basicstyle=\scriptsize\ttfamily,
  morestring=[s]{"}{"},
  morecomment=[s]{?}{?},
  morecomment=[s]{!--}{--},
  commentstyle=\color{darkgreen},
  moredelim=[s][\color{black}]{>}{<},
  moredelim=[s][\color{red}]{\ }{=},
  stringstyle=\color{maroon},
  identifierstyle=\color{blue}
}

\lstdefinelanguage{json}
{
    basicstyle=\scriptsize\ttfamily,
    tabsize=2,
    %numbers=left,
    %numberstyle=\scriptsize,
    %stepnumber=1,
    %numbersep=8pt,
    showstringspaces=false,
    breaklines=true,
    frame=lines,
    literate=
     *{0}{{{\color{numb}0}}}{1}
      {1}{{{\color{numb}1}}}{1}
      {2}{{{\color{numb}2}}}{1}
      {3}{{{\color{numb}3}}}{1}
      {4}{{{\color{numb}4}}}{1}
      {5}{{{\color{numb}5}}}{1}
      {6}{{{\color{numb}6}}}{1}
      {7}{{{\color{numb}7}}}{1}
      {8}{{{\color{numb}8}}}{1}
      {9}{{{\color{numb}9}}}{1}
      {:}{{{\color{punct}{:}}}}{1}
      {,}{{{\color{punct}{,}}}}{1}
      {\{}{{{\color{delim}{\{}}}}{1}
      {\}}{{{\color{delim}{\}}}}}{1}
      {[}{{{\color{delim}{[}}}}{1}
      {]}{{{\color{delim}{]}}}}{1},
}

\lstdefinelanguage{JavaScript}{
  keywords={typeof, new, true, false, catch, function, return, null, catch, switch, var, if, in, while, do, else, case, break},
  keywordstyle=\color{blue}\bfseries,
  ndkeywords={class, export, boolean, throw, implements, import, this},
  ndkeywordstyle=\color{darkgray}\bfseries,
  identifierstyle=\color{black},
  sensitive=false,
  comment=[l]{//},
  morecomment=[s]{/*}{*/},
  commentstyle=\color{green}\ttfamily,
  stringstyle=\color{red}\ttfamily,
  morestring=[b]',
  morestring=[b]"
}

%% this would set javascript as default as well as other params
%\lstset{
% language=JavaScript,
% backgroundcolor=\color{lightgray},
% extendedchars=true,
% basicstyle=\footnotesize\ttfamily,
% showstringspaces=false,
% showspaces=false,
% numbers=left,
% numberstyle=\footnotesize,
% numbersep=9pt,
% tabsize=2,
% breaklines=true,
% showtabs=false,
% captionpos=b
%}


\hypersetup{urlcolor=blue, colorlinks=true} % Colors hyperlinks in blue - change to black if annoying
\title{\ttitle} % Defines the thesis title - don't touch this

\begin{document}

\frontmatter % Use roman page numbering style (i, ii, iii, iv...) for the pre-content pages

\setstretch{1.3} % Line spacing of 1.3

% Define the page headers using the FancyHdr package and set up for one-sided printing
\fancyhead{} % Clears all page headers and footers
\rhead{\thepage} % Sets the right side header to show the page number
\lhead{} % Clears the left side page header

\pagestyle{fancy} % Finally, use the "fancy" page style to implement the FancyHdr headers

\newcommand{\HRule}{\rule{\linewidth}{0.5mm}} % New command to make the lines in the title page

% PDF meta-data
\hypersetup{pdftitle={\ttitle}}
\hypersetup{pdfsubject=\subjectname}
\hypersetup{pdfauthor=\authornames}
\hypersetup{pdfkeywords=\keywordnames}

%-----------------
%	TITLE PAGE
%---------------

\begin{titlepage}
\begin{center}

\textsc{\Large \univname}\\[1cm] % University name

\HRule \\[0.4cm] % Horizontal line
{\huge \bfseries \ttitle}\\[0.4cm] % Thesis title
\HRule \\[1.5cm] % Horizontal line
 
\begin{minipage}{0.4\textwidth}
\begin{flushleft} \large
\emph{Authors:}\\
\authornames % Author name - add \href bracket with link to add a website
\end{flushleft}
\end{minipage}
\begin{minipage}{0.4\textwidth}
\begin{flushright} \large
\emph{Supervisor:} \\
\supname % Supervisor name - add \href bracket with link to add a website
\end{flushright}
\end{minipage}\\[11cm]
\groupname\\
 
{\large November 2013}\\[1cm] % Date

\begin{tabular}{lr}
   \includegraphics[scale=0.6]{./Figures/NTNU-logo.pdf} 
   \hspace{3 cm} &
   \includegraphics[scale=0.9]{./Figures/HDIR-LOGO-HORISONTAL-ENG-CMYK.pdf} \\
\end{tabular}
\end{center}
\end{titlepage}
\clearpage % Start a new page

%---------------
%	ABSTRACT PAGE
%----------------

\addtotoc{Abstract} % Add the "Abstract" page entry to the Contents

\abstract{\addtocontents{toc}{\vspace{1em}}

This report is the final product for the course TDT4290 Customer driven project. 
Our customer, The Directorate of Health, Department of the Health Portal, gave us the task of planning, designing, describing and developing a prototype of the National Integration Platform for Citizen Centric eHealth in Norway. 
The intention of such a platform is to enable citizens’ the ability to publish and fetch information they produce into the government run citizen centric health portal helsenorge.no. 
By creating this platform The Directorate of Health will digitize the Norwegian health care which is one of their main goals. 

Recent years a movement called quantified self has started emerging \cite{quantified} leveraging that citizens has smartphones and tablets readily available that are able to capture and store health data. 
This data together with health information gathered by educated professionals is what The Directorate of Health wants to collect at the eHealth portal.
However national IT projects usually have a high degree of complexity, security and scaling needs which are central to the design of the platform that makes projects like this difficult. 

Our customer specified that there are two types of stakeholders who are central for the pursuit of this project and that should be the goal of our prototype:
1. Educated health professionals
2. Developers.
The prototype should have a demonstration that is easy for the first group to grasp and also make it appealing and technical for the second group.

The motivation for working with this project is to innovate and create a new type of platform that can be used to better analyze citizens health. 
This report examines multiple similar or related solutions for publishing and fetching health data.
Our proposed solution consists of a web application programming interface (API) to transfer JSON data models representing health data such as heart rate or weight.
We produced a number of Android applications and a web portal to study the feasibility and practical use of the portal.
We also integrated our solution with Microsoft HealthVault.
Our results outline the viability of the chosen approach which will require further work to be suitable for a national scale project.
The next step for this project should be to analyse the value of the type of data that can be collected in this platform. 

\clearpage % Start a new page

%------------------
%	PreFace
%------------------

\addtotoc{Preface} 

\preface{\addtocontents{toc}{\vspace{1em}}
This report is the final product for the course TDT4290 Customer driven project at NTNU. 
The project was executed on behalf of The Directorate of Health, Department of the Health Portal.

The group consisted of three students from NTNU.
The team would like to thank Meng Zhu for his advices and supervision.
We would also like to thank the customer representative Mr. Helge T. Blindheim for his participation and great guidance.
}
\clearpage % Start a new page

%-----------------------------------------
%	LIST OF CONTENTS/FIGURES/TABLES PAGES
%------------------------------------------

\pagestyle{fancy} % The page style headers have been "empty" all this time, now use the "fancy" headers as defined before to bring them back

\addtotoc{Contents}
\addtocontents{toc}{\vspace{1em}}
\lhead{\emph{Contents}} % Set the left side page header to "Contents"
\tableofcontents % Write out the Table of Contents

\lhead{\emph{List of Figures}} % Set the left side page header to "List of Figures"
\listoffigures % Write out the List of Figures

\lhead{\emph{List of Tables}} % Set the left side page header to "List of Tables"
\listoftables % Write out the List of Tables

%----------------
%	ABBREVIATIONS
%-------------------

\clearpage % Start a new page

\setstretch{1.5} % Set the line spacing to 1.5, this makes the following tables easier to read

\lhead{\emph{Abbreviations}} % Set the left side page header to "Abbreviations"
\listofsymbols{ll} % Include a list of Abbreviations (a table of two columns)
{
\textbf{AJAX} & \textbf{A}synchronous \textbf{Ja}vaScript and \textbf{X}ML \\
\textbf{API} & \textbf{A}pplication \textbf{P}rogramming \textbf{I}nterface \\
\textbf{App} & \textbf{A}pplication \\
\textbf{BPM} & \textbf{B}eats \textbf{P}er \textbf{M}inute \\
\textbf{CSS} & \textbf{C}ascading \textbf{S}tyle \textbf{S}heets \\
\textbf{DOM} & \textbf{D}ocument \textbf{O}bject \textbf{M}odel \\
\textbf{HTML} & \textbf{H}yper\textbf{T}ext \textbf{M}arkup \textbf{L}anguage \\
\textbf{HV} &  Microsoft's \textbf{H}ealth\textbf{V}ault. \\
\textbf{IDI} & \textbf{I}nstitutt for \textbf{D}atateknikk og \textbf{I}nformasjonsvitenskap  \\
\textbf{IP} & \textbf{I}ntegration \textbf{P}latform \\
\textbf{JSP} & \textbf{J}ava\textbf{S}erver \textbf{P}ages \\
\textbf{MVC} & \textbf{M}odel-\textbf{v}iew-\textbf{c}ontroller \\
% Should we include the english name? Department of Computer and Information Science \\
\textbf{NIP} & \textbf{N}ational \textbf{I}ntegration \textbf{P}latform \\
\textbf{NIPEN} & \textbf{N}ational \textbf{I}ntegration \textbf{P}latform for \textbf{e}Health in \textbf{N}orway \\
\textbf{NTNU} & \textbf{N}orges \textbf{T}eknisk-\textbf{N}aturvitenskalige \textbf{U}niversitet \\
% Should we include the english name? \\Norwegian University of Science and Technology\\
\textbf{POM} & \textbf{P}roject \textbf{O}bject \textbf{M}odel \\
\textbf{REST} & \textbf{R}epresentational \textbf{S}tate \textbf{T}ransfer \\
\textbf{SDK} & \textbf{S}oftware \textbf{D}evelopment \textbf{K}it \\
\textbf{SSL} & \textbf{S}ecure \textbf{S}ockets \textbf{L}ayer \\
\textbf{TLS} & \textbf{T}ransport \textbf{L}ayer \textbf{S}ecurity \\
\textbf{URL} & \textbf{U}niform \textbf{R}esource \textbf{L}ocator \\
}

%-----------------------------
%	THESIS CONTENT - CHAPTERS
%------------------------------

\mainmatter % Begin numeric (1,2,3...) page numbering

\pagestyle{fancy} % Return the page headers back to the "fancy" style

% Include the chapters of the thesis as separate files from the Chapters folder
% Uncomment the lines as you write the chapters


\chapter{Introduction} 
\label{Introduction} 

\lhead{Chapter 1. \emph{Introduction}}

This chapter provides a description of the project, its goals, duration, background and involved parties.
At the end, we present a brief description of every other chapter in this document.

%-------------------------
\section{Project description}
\label{section:description}

The purpose of the project was to design, develop and document an integration platform for citizen eHealth\cite{ehealth}. 
The aim of such platform is to integrate digital information about citizens' health \iffalse which is\fi collected using popular
devices such as smartphones and tablets.
%citizens to publish data related to their health produced by devices or third party solutions. 
%It is very common today for people to log health data from mobile phones and tablets in their possession.
The project description as provided by the customer can be found in Appendix \ref{AppendixA}.

%% very relevant but to be written elsewhere. (project goal)
%For the client it was important that this project produced a prototype that could be demonstrated to:
%1. Educated health professionals
%2. Developers

%% very relevant but to be written elsewhere. (project limitations?)
% Security is of high importance when dealing with citizen health data. This was not made a requirement of the project beacue the assignment had to be scaled down because the group only consisted of three members. We will however discuss how you could add a secure layer to the working solution. 

The name we chose for the project is \textbf{NIPEN}.\newline
It is simply an acronym for \textbf{N}ational \textbf{I}ntegration \textbf{P}latform for \textbf{e}Health
in \textbf{N}orway.

%-------------------------
\section{Project customer}
\label{section:client}

The project customer was the Department of the Health portal, Norwegian Directorate of Health (Helsedirektoratet)
located in Oslo, Norway. The Directorate has, among other, the task of digitalizing Norway's health care system
by providing services for both specialists and citizens. The customer was represented by Mr. Helge T. Blindheim.
His contact information is shown in table \ref{table:contact}.

% --- relevant but to be detailed in another chapter (proj. mgmt) - emanuele
%This infered that our weekly meeting had to be heald over teleconferencing.

%% COMMENTED OUT
\iffalse
\begin{table}[h]
\begin{center}
\begin{tabular}{ | l | l | l |  l | }
  \hline
  Name & Phone & E-mail & \\
  \hline\noalign{\smallskip}\hline
  Helge T. Blindheim	& 46675321 & Helge.T.Blindheim@helsedir.no \\
  \hline
\end{tabular}
\end{center}
\caption{Customer representative}
\label{table:client}
\end{table}
\fi

\section{Involved parties}
\label{section:parties}

The people involved in this project were the customer, the team and the supervisor.
The customer, introduced in section \ref{section:client}, was represented by Mr. Helge T. Blindheim.
The team consisted of three students from the Department of Computer and Information Science (IDI) at the Norwegian University of Science and Technology (NTNU). The group was supervised by PhD. candidate Zhu Meng.
See table \ref{table:contact} for contact information.

\begin{table}[h]
\begin{center}
\begin{tabular}{ | l | l | l | p{2.5cm} | }
  \hline
  Name & Phone & E-mail & Party \\
  \hline\noalign{\smallskip}\hline
  Helge T. Blindheim	& 46675321		& Helge.T.Blindheim@helsedir.no		& Customer\\
  Meng Zhu				& 73551189		& zhumeng@idi.ntnu.no				& Advisor\\
  Anders Olsen Sandvik	& 91824583		& andsan@stud.ntnu.no				& Team member \\
  Emanuele Di Santo		& none			& lemrey@gmail.com					& Team member \\
  Sebastian Zalewski	& 95107928		& zalewski@stud.ntnu.no				& Team member \\
  \hline
\end{tabular}
\end{center}
\caption{Contact information}
\label{table:contact}
\end{table}

%% COMMENTED OUT
\iffalse
\begin{table}[h]
\begin{center}
\begin{tabular}{ | l | l | l | }
  \hline
  Name & Phone & E-mail \\
  \hline\noalign{\smallskip}\hline
  Zhu Meng	& 73551189 & zhumeng@idi.ntnu.no \\
  \hline
\end{tabular}
\end{center}
\caption{Student advisor}
\label{table:advisor}
\end{table}
\fi


%--------------
\section{Project background}
\label{section:background}

Digital healthcare is about using information technologies to provide solutions to problems in healthcare.
This definition includes a lot of different domains among which is eHealth:
\begin{quote}
an emerging field in the intersection of medical informatics, public health and business, referring to health services
and information delivered or enhanced through the Internet and related technologies.\citep{ehealth}
\end{quote}
eHealth projects, especially on a national scale, are therefore long, complex and inherently costly.
At the same time the progress in information technology has made available powerful and yet cheap devices which
can be used to monitor people's health. These devices are nowdays widespread and used by a large part of the
population. The Department of Health Portal (our customer) is interested in investigating the
possibilities and advantages of using these devices in national eHealth projects.

%This project is part of the Customer Drivent Project (TDT4290 2013) at NTNU.

\section{Project objective}
\label{section:objective}

The project had multiple goals.
The purpose of the course which this project is part of is to let students acquire practical experience in development
of a medium-large software project, including experience in project management, group dynamics and customer relations.

The aim of the customer was to investigate the possibilities for national eHealth projects to use widely available
devices such as smartphones and tables. The customer also requested a prototype system that could demonstrate such
possibilities and stimulate further work.

Expected deliverables:
\begin{itemize}
\item an extensive documentation of the process and the product, called 'Project report'
\item an Integration platform for health information
\item a number of prototype applications using the Integration platform
\end{itemize}

\section{Duration}
\label{section:duration}
The project started on August the 21th. The final presentation will be held on November the 21th.
That makes a total of about 13 weeks to work on the project.
We were advised to dedicate at least 24 hours per week to the project by our teachers.
That makes a total of 312 hours per student. Because we are a group of three students we estimate
the project to take at least 936 hours of work to be completed.

\begin{itemize}
\item Start date: 21.08.2013
\item End date: 21.11.2013
\item Total hours: 936
\item Months: 3
\item Days: 93
\end{itemize}

\section{Summary}

See below for a summary of each chapter in this document.

\begin{itemize}
	\item \textbf{Project management}\newline
		This chapter contains details about the project plan and organization.
	\item \textbf{Preliminary studies}\newline
		This chapter covers our preliminary studies, including technologies
		and similar solutions. At the end of each section we draw a conclusion
		on each topic covered.
	\item \textbf{Requirements and specification}\newline
		The purpose of this chapter is to formally identify the project stakeholders
		and requirements, both functional and non-functional.
	\item \textbf{System architecture}\newline
		Provides an insight into the architecture of the system,
		detailing every part of it and discussing the rationale behind its design.
	\item \textbf{Security}\newline
		The aim of this chapter is to describe some security practices that
		should be put in place for this type of projects.
	\item \textbf{Sprint 0-5}\newline
		Describe in detail the goals, planning and results of each sprint.
		Every chapter is concluded by an internal evaluation of the sprint.
	\item \textbf{Conclusion and further work}\newline
		In this chapter we state our results and draw a conclusion.
		We point out at the limitations of our product and possible further work
		that can be done to improve it.
	\item \textbf{Evaluation}\newline
		This chapter summarizes our evaluation about every aspect of this project:
		the results, the learning outcome, the social dynamics\ldots
\end{itemize}
\chapter{Project management}

\label{ch:management}
\lhead{Chapter 2. \emph{Project management}}

This chapter describes the planning of the project. Section \ref{section:planning} is about project planning,
section \ref{section:organization} covers project organizazion.

\section{Project planning}
\label{section:planning}
This section describes the planning of the project including schedule, resources and tool selection.

\subsection{Project goals}
The project has various goals. One is the delivery of a report that describes the all phases and relevant aspects regarding the development of the project. Another one is the development of the integration platform together with two or three prototype applications which can be used to demonstrate the functionality of the system and possibily stimulate further development of a national integration platform for eHealth.

Expected deliverables:

\begin{itemize}
\item Project report
\item Integration platform
\item Two or three prototype applications
\end{itemize}


%\subsection{Limitations}

\subsection{Project schedule}

\textbf{Milestones} \newline
We have identified four milestones for the project. These are associated with important events in the project lifetime and can help us to better monitor its progress. Project's milestones are shown in table \ref{table:milestones}.

\begin{table}[h]
\begin{center}
\begin{tabular}{ | l | l | l | }
  \hline
  Milestone & Date & Description \\
  \hline\noalign{\smallskip}\noalign{\smallskip}\hline
  M1 & 20.09.2013 & Basic implementation of the platform and one prototype. \\ 
  M2 & 08.10.2013 & Mid-term report delivery. \\
  M3 & 04.11.2013 & Feature freeze. \\
  M4 & 21.11.2013 & Project demonstration. Final report delivery. \\
  \hline
\end{tabular}
\end{center}
\caption{Project milestones}
\label{table:milestones}
\end{table}

\textbf{Sprints schedule} \newline
Table \ref{table:sprints} contains an overview of the Sprints in the project.

\begin{table}[h]
\begin{center}
\begin{tabular}{ | l | l | l | }
  \hline
  Sprint & Start date & End date \\
  \hline\noalign{\smallskip}\noalign{\smallskip}\hline
  0 &  &  \\ 
  1 & 09.09.2013 & 22.09.2013 \\
  2 & 23.09.2013 & 06.10.2013 \\
  3 & 07.09.2013 & 20.10.2013 \\
  4 & 21.10.2013 & 04.11.2013 \\
  5 & 05.11.2013 & 21.11.2013 \\
  \hline
\end{tabular}
\end{center}
\caption{Planned sprints}
\label{table:sprints}
\end{table}

\subsection{Work plan}
\newpage
\begin{landscape}
\begin{figure}[h]
\includegraphics[scale=0.66]{../Figures/gantt-diagram.png}
\caption{Gantt diagram}
\label{figure:work-splan}
\end{figure}
\end{landscape}

\subsection{Resources}
The team consists of three members from different educational backgrounds and with different skills.
We have thus created a skill matrix regarding the technologies involved in the project in order to better distribute the work among us.
The skills are valued on a scale of zero to five based on criterias according to table \ref{table:skillscale}.
Each value corresponds to a certain degree of proficiency, detailed in table \ref{table:proficiency}.
The skill matrix is shown in table \ref{table:skillmatrix}.

\begin{table}[h]
\begin{center}
\begin{tabular}{ | c | l | l | }
  \hline
  Level & Proficiency & Criteria \\
  \hline\noalign{\smallskip}\noalign{\smallskip}\hline
  0 & None		& Has no clue what we're talking about. \\
  1 & Basic		& Has read about it.\\
  2 & Fair		& Has studied it at school.\\
  3 & Good		& Personal interest, use in small-medium sized projects.\\
  4 & Very good	& Strong personal interest, use in medium-large sized projects. \\
  5 & Excellent	& 10+ Years of work experience. \\
  \hline
\end{tabular}
\end{center}
\caption{Skill's scale explanation}
\label{table:skillscale}
\end{table}

\begin{table}[h]
\begin{center}
\begin{tabular}{ | l | l | }
  \hline
  Proficiency & Description \\
  \hline\noalign{\smallskip}\noalign{\smallskip}\hline
  None		& Cannot perform the task \\
  Basic		& Can perform the task with some help from other people \\
  Fair		& Can perform the task almost indipendently \\
  Good		& Can perform the task without any help \\
  Very good	& Can help others to complete the task \\
  Excellent	& Can help others to complete the task \\
  \hline
\end{tabular}
\end{center}
\caption{Proficiency level description}
\label{table:proficiency}
\end{table}

\begin{table}[h]
\begin{center}
\begin{tabular}{ | l | c | c | c | c | c | c | c | c | }
  \hline
  Team member & Java & JS & Android & Database & CSS & AJAX & Spring & LaTeX \\
  \hline\noalign{\smallskip}\noalign{\smallskip}\hline
  Anders & 4 & 4 & 3 & 4 & 4 & 2 & 3 & 1 \\
  Emanuele & 4 & 1 & 3 & 2 & 1 & 1 & 0 & 3 \\
  Sebastian & ? & ? & ? & ? & ? & ? & ? & ? \\
  \hline
\end{tabular}
\end{center}
\caption{Skill matrix}
\label{table:skillmatrix}
\end{table}

\subsection{Tool selection}
This section will describe the different tools we used throught the project.


\begin{description}

\item[Git and GitHub]
Git is one of the most used version control systems. Although it is mostly used for code, it can be used to manage other types of file such as LaTeX. GitHub is a repository hosting service for Git projects which offers a web-based fronted. We have decided to use Git because of the familiarity we had with it. We used Git for managing both the code and the report.

\item[Sublime Text]
Sublime Text is a popular, cross-platform editor which offers advanced features.
We used Sublime Text for editing the report and for coding JavaScript, CSS and AJAX.

\item[IntelliJ IDEA]
A popular, feature-rich Java IDE. We used it to develop the Android applications and for Spring coding.

\item[Apache Maven]
Maven is an open and cross-platform build system for Java projects which, among other things, takes care of dependecies. Maven is powerful and well documented and seemed an optimal choice for our project also because of the familiarity we had with it.

\item[LaTeX suite]
LaTeX is a powerful language used to prepare documents which is widely used in academical environments. LaTeX takes care of the formatting of the document allowing the writer to focus on the contents, moreover LaTeX documents can be written with any text editor.
We chose LaTeX because of the familiarity we had with it.
%LaTeX files can then be compiled to PDF files.

\item[Google Drive and Google Docs]
Google Drive is a service which lets you access your files from everywhere and share them with other people. It is tightly integrated with Google Docs, a web-based productivity suite. We used Google Drive to share documents within the team and with the customer and supervisor.
Many of the documents and diagrams produced for the project were created using Google Docs.

\item[Skype]
A popular instant-messaging client that also supports audio and video communication.
Skype was used throught the project to keep in touch with the customer and team members.

\item[Travis CI]
Travis CI is an automated build service...
% we should write two more lines here.

\item[Facebook]
Facebook is a social networking service. The team created a group on Facebook that was used as a bulleting board for internal communications.

%\item[Violet UML] A free UML editor.
%Balsamiq Mockups looks great, it would be cool to do some mockup with this!!
%Lucidchart maybe we can opt for a free solution like violet UML.

\end{description}


\section{Project organization}
\label{section:organization}
This section describes the organizational aspects of the project such as roles description, allocation and meetings schedule.

\subsection{Role description}
We have identified a number of roles in this project. Although each role has different, distinct responsabilities, we see all roles as equally important for the success of this project.

\begin{description}
\item[Project manager]
\item[Scrum master]
\item[Product owner]
%\item[Document responsible]
%\item[System architect]
%\item[]
\item[Secretary]
Responsible for taking notes during meetins and booking rooms.
\item[Quality assurance]
Ensures that quality practices are in place and use.
\item[Web developer]
Responsible for web development.
\item[Droid developer]
Responsible for Android development.
\item[Supervisor relations]
Responsible for communication with the supervisor.
\item[Customer relations]
Responsible for contacting the customer.
\end{description}

\subsection{Role allocation}
Roles were allocated mostly basing on team members' competences.
Those roles that were left unassigned were taken up by members who volunteered for them upon approval of other team members.
Because we were only three people, all of us had more than one role. Sometimes, roles were shared.

\begin{table}
\begin{center}
\begin{tabular}{ | l | l | }
  \hline
  Member & Roles \\
  \hline\noalign{\smallskip}\noalign{\smallskip}\hline
  Anders Olsen Sandvik  &  \\
  Emanuele Di Santo     &  \\
  Sebastian Zalewski    &  \\
  \hline
\end{tabular}
\end{center}
\caption{Members' roles}
\label{table:roles}
\end{table}

\subsection{Weekly schedule}
During the early stages of the project the team has agreed on a meeting schedule.

\begin{description}
\item[Monday] \hfill from 10 to 19 \\
\textbf{Weekly startup meeting} from 1000 to 1100.\newline
At the begin of each week we would review the progress done during the previous and discuss the plan for the next.
We also discussed any problems encountered.
This internal meeting would follow the last Status Report document.
\textbf{Customer meeting} from 1100 to 1200.\newline
During the meeting with customer we would discuss with him the work done and share our plan for the next week.
This was done in order to help us prioritize our work to better suit the customer's interest.
\textbf{Supervisor meeting} from 1200 to 1300.
\item[Wednesday] \hfill \\
\textbf{Collective work session} from 15 to 1900. (4 hours)
\end{description}


\section{Quality assurance}
\subsection{Templates}
\subsection{Customer relations}
\subsection{Supervisor relations}

\section{Risk management}

The risk analysis is good but we can improve it with some strategy to handle the risk once it has happened.

\chapter{Preliminary Studies}
\label{Preliminary Studies}
\lhead{Chapter 5. \emph{Preliminary Studies}}

This chapter contains

3.1 dev metho
3.1.1 waterfall
3.1.2 scrum
3.2 existing solutions
3.2.1 HealthVault
3.2.2 open ehealth
3.2.3 human api
3.3 tech
3.3.1 server
3.3.2 database
3.3.3 web page
3.3.4 android
3.4 testing


%-----------
% SECTION 1
%-----------

\section{Development Methodology}

TODO

%-----------------------------------
%	SUBSECTION 1
%-----------------------------------
\subsection{Waterfall Model}

The waterfall model is a software development process where each task is performed in a sequential order.
Before moving to the next phase the preceding task needs to be finished.
The progress of the project is seen as flowing downwards through the different phases, hence the name waterfall.
In the original model the phases consisted of seven different tasks:

\begin{enumerate}
\item Requirements specification
\item Design
\item Construction (implementation or coding)
\item Integration
\item Testing and debugging
\item Installation
\item Maintenance
\end{enumerate}

Because each phase needs to be perfected and completed before moving to the next phase, this brings up some difficulties if the requirements were to change during the development process. 
However the model is easily understandable, structured, and disciplined. 
All the phases are divided into different sections, and this makes it easier to understand the progress of the project.
In practice it can be very hard to adapt to this kind of development model. 
It can be hard for a system designer to predict future implementation difficulties of a type of design, hence the design of the system may change during the process.
Another problem is that the customer is not always sure about the system requirements, and often will the customer change them during the development.

%-----------------------------------
%	SUBSECTION 2
%-----------------------------------

\subsection{SCRUM Model}

TODO

%-----------------------------------
%	SUBSECTION 3
%-----------------------------------
\subsection{Conclusion}

TODO

%----------------------------------------------------------------------------------------
%	SECTION 2
%----------------------------------------------------------------------------------------

\section{Existing Solutions}

TODO

%-----------------------------------
%	SUBSECTION 1
%-----------------------------------
\subsection{HealthVault}

TODO

%-----------------------------------
%	SUBSECTION 2
%-----------------------------------
\subsection{Open eHealth Foundation}

TODO

%-----------------------------------
%	SUBSECTION 3
%-----------------------------------
\subsection{human/api}

The human API is a platform for human health data. 
They have an API that contains multiple different well defined JSON strings for different kinds of human related data.
Each JSON string contains all the necessary information that is needed to represent each type of health data.
For example heart rate is defined by an id, user id, time, value and unit in the following way:

\begin{verbatim}
{
  "id": "string",
  "userId": "string",
  "time": "date",
  "value": "int",
  "unit": "string"
}
\end{verbatim}

%-----------------------------------
%	SUBSECTION 4
%-----------------------------------
\subsection{Conclusion}

TODO

%----------------------------------------------------------------------------------------
%	SECTION 3
%----------------------------------------------------------------------------------------

\section{Technologies}

TODO

%-----------------------------------
%	SUBSECTION 1
%-----------------------------------
\subsection{Server}

Java \\
Spring Framework \\
TODO

Apache Tomcat Server \\ 
TODO

%-----------------------------------
%	SUBSECTION 3
%-----------------------------------
\subsection{Database}

\textbf{MySQL}

TODO

%-----------------------------------
%	SUBSECTION 4
%-----------------------------------
\subsection{Web Page}

\textbf{HTML5}

HTML is the standard World Wide Web's markup language.
It is used to structure and visualize web pages on the internet.

\textbf{CSS3}

CSS describes the look and format of a document written in HTML.

\textbf{Javascript}

TODO

\textbf{JQuery}

TODO

\textbf{Chart.js}

Chart.js is a Javascript library for creating graphs TODO

%-----------------------------------
%	SUBSECTION 5
%-----------------------------------
\subsection{Mobile Technologies}

Android SDK \\
TODO

%-----------------------------------
%	SUBSECTION 6
%-----------------------------------
\subsection{Conclusion}

TODO

%----------------------------------------------------------------------------------------
%	SECTION 4
%----------------------------------------------------------------------------------------

\section{Testing}

TODO

%-----------------------------------
%	SUBSECTION 2
%-----------------------------------
\subsection{Conclusion}

%----------------------------------------------------------------------------------------
%	SECTION 5
%----------------------------------------------------------------------------------------

\section{Summary}

TODO
 
\chapter{Requirements specification}

\label{ch:requirements}
\lhead{Chapter 4. \emph{Requirements specification}}

This chapter describes the requirements for the product. These include both functional and non-functional requirements.
Functional requirements can be found in section \ref{section:functionalreq},
non-functional requirements in section \ref{section:nonfunctionalreq}.

Throughout this chapter both requirement's priority and complexity have a textual description which can be
\begin{itemize}
\item High
\item Medium (abbrev. Med)
\item Low
\end{itemize}

The use of such terms is described in table \ref{table:priorities} (priorities) and table \ref{table:complexity} (complexity).

\begin{table}[H]
\begin{center}
\begin{tabular}{ | c | p{12.5cm} | }
  \hline
  Priority & Description \\
  \hline\noalign{\smallskip}\noalign{\smallskip}\hline
  High & An essential requirement. The product \textbf{must fulfill} the requirement in order to be satisfactory. \\
  Medium & A useful requirement. The product \textbf{should fulfill} the requirement to maximise effectiveness. \\
  Low & A desiderabe requirement. The product \textbf{could fulfill} the requirement to be more interesting for certain stakeholders. \\
  \hline
\end{tabular}
\end{center}
\caption{Priority descriptions}
\label{table:priorities}
\end{table}

\begin{table}[H]
\begin{center}
\begin{tabular}{ | c | p{12.5cm} | }
  \hline
  Complexity & Description \\
  \hline\noalign{\smallskip}\noalign{\smallskip}\hline
  High & The requirement is difficult to implement. It will take a considerable amount of time. (more than 30 hours?) \\
  Medium & The requirement is moderately difficult. It will take some time. (from 10 to 30 hours?) \\
  Low & The requirement is easy and can be achieved in a short amount of time. (10h or less?) \\
  \hline
\end{tabular}
\end{center}
\caption{Complexity descriptions}
\label{table:complexity}
\end{table}

%-----

\section{Stakeholders}
\label{section:stakeholders}

This section contains the stakeholders of our system.
A short description of the role of each stakeholder is given, and what concerns they might have.

\subsection{Customer}
The customer is the one that is going to guide us to the solution he wants.
His concern is that we should maintain an effective and good communication with him, so he understands our progress.
We should also document the advancement we are making and have a clear system architecture for him.
His main concern is to get a working prototype of our system according to his requirements.

\subsection{Course evaluators}
The course evaluators are the one that are going to evaluate our project. 
Their concern is that the project should be well documented and understandable, and that all the deliveries should be on schedule.
Good communication is also of importance.

\subsection{Developers}
The developers include the creators of the system, and third party developers.
The concern of the creators of the system is to have some requirements and guidance from the customer.
In this way the goal of the project will be clear.
The third party developers are going to integrate our application to their system.
Their concern is that our integration platform should easily be connected with their applications.

\subsection{Users}
The users are the one that are going to use the front-end, i.e. the web page that visualizes the data from the integration platform.
Their concern is that their data should be available for them, and visualized in an intuitive and comprehensible way.

%-----

\section{Funcional requirements}
\label{section:functionalreq}

This section describes the functional requirements for the product.
Each requirements has a priority which helps us identify the focus of our work.
Priorities can be reviewed based on customer's feedback. We tried to prioritize those requirements corresponding to the functionality of the system that the customer expressed most interest about. In order to better organize our workload we also assigned a \'difficulty\' to each requirement.

\textbf{Functional requirements for Integration platform}

\begin{table}[H]
\begin{center}
\begin{tabular}{ | c | p{9cm} | c | c | }
  \hline
  ID & Description & Difficulty & Priority\\
  \hline\noalign{\smallskip}\noalign{\smallskip}\hline
  FIP1	& The IP shall support REST endpoints for receiving heart rate data models expressed as JSON strings	& Med	& High \\
  FIP2	& The IP shall support REST endpoints for receiving weight data models expressed as JSON strings 		& High	& High \\
  FIP3	& The IP shall support REST endpoints for forwarding heart rate models (using JSON) to other systems.	& Med	& High \\
  FPI4	& The IP shall support REST endpoints for forwarding weight models (using JSON).						& Med	& High \\
  \hline
\end{tabular}
\end{center}
\caption{Functional requirements for Integration platform}
\label{table:reqip}
\end{table}

\textbf{Functional requirements for the web frontend}

\begin{table}[H]
\begin{center}
\begin{tabular}{ | c | p{9cm} | c | c |}
  \hline
  ID & Description & Difficulty & Priority\\
  \hline\noalign{\smallskip}\noalign{\smallskip}\hline
  FW1	& The web frontend shall display the data stored by the Integration Platform using charts.	& Med	& High \\
  FW2	& The web frontend shall use Helsenorge color palette										& Low	& Low \\
  \hline
\end{tabular}
\end{center}
\caption{Functional requirements for the web frontend}
\label{table:reqfrontend}
\end{table}

\textbf{Functional requirements for Heart rate application}

\begin{table}[H]
\begin{center}
\begin{tabular}{ | c | p{9cm} | c | c |}
  \hline
  ID & Description & Difficulty & Priority\\
  \hline\noalign{\smallskip}\noalign{\smallskip}\hline
  FHR1	& The application shall measure user’s heart rate using the device camera	& Med	& High \\
  FHR2	& The application shall display acquired measurements on the screen.		& Low	& Low \\
  FHR3	& The application shall forward the data to the IP using its REST endpoint. & Med	& High \\
  \hline
\end{tabular}
\end{center}
\caption{Functional requirements for Heart rate application}
\label{table:reqheartrate}
\end{table}

\textbf{Functional requirements for Weight application}

\begin{table}[H]
\begin{center}
\begin{tabular}{ | c | p{9cm} | c | c |}
  \hline
  ID & Description & Difficulty & Priority\\
  \hline\noalign{\smallskip}\noalign{\smallskip}\hline
  FHV1	& The application shall fetch weight data from HealthVault.						      & Med	& High \\
  FHV2	& The application shall show the user the data it has fetched.              & Low	& Med \\
  FHV3	& The application shall forward the data to the IP using its REST endpoint. & Med	& High \\
  \hline
\end{tabular}
\end{center}
\caption{Functional requirements for Weight application}
\label{table:reqweight}
\end{table}

%-----

\section{Non-funcional requirements}
\label{section:nonfunctionalreq}

This section outlines non-functional requirements (quality attributes) for the product.
In order to provide a better overview, we have organized them in categories.

\begin{table}[H]
\begin{center}
\begin{tabular}{ | c | c |p{6.5cm} | c | c |}
  \hline
  ID & Category & Description & Difficulty & Priority\\
  \hline\noalign{\smallskip}\noalign{\smallskip}\hline
  NF1 & Documentation & The system shall be thoroughly documented, both at the code level and by the document ‘project report’.
  & High & High \\
  NF2 & Documentation & Although security and privacy are not requirements for the product they are important topics to be discussed in the documentation.
  & Med & High \\
  NF3 & Open-source	& The product shall be released under a permissive license approved by the product owner.
  & Low & High \\
  NF4 & Interoperability & The system shall provide a good degree of interoperability. Third party application developers should be put in the condition to develop third party (interoperable) solutions rapidly.
  & Med & High \\
  NF5 & Interoperability & A number of two-three prototype applications shall be developed in order to showcase the functionality of the system.
  & Med & High \\
  NF6 & Accessibility & The web-frontend should have a good degree of accessibility. It should have a rather simple design and use a user-provided palette.
  & Low & Low \\
  \hline
\end{tabular}
\end{center}
\caption{Non-functional requirements}
\label{table:reqnonfunc}
\end{table}

% -----

\section{Requirements Validation}
(Sommerville 7.3)
How to ensure that a Requirements Document does define 
the customers’ requirements (check against desirable 
features)
• Comprehensibility - do customers understand the 
requirements?
• Verifiability (testability) - can requirements be 
realistically tested? If you cannot test it, it is not a 
requirement!
• Traceability - source of the requirement? Important 
when making changes.
• Adaptability (changeability) - what is the impact of 
changing a particular requirement?

%-----

\section{Use cases}

In this section we are going through the different use cases of our applications.
The section consists of use cases for the front-end, heart rate application and the weight application.

\subsection{Front-end}

The use case diagram for the front-end is shown in figure \ref{figure:use-case-diagram-front-end}.
From the use case diagram we see that the front-end has two main objectives.
The first one is to fetch data from our integration platform, and the second one is to visualize this data for the user.
In table \ref{table:use-case-view-data} a textual use case is given for the front-end.

\begin{figure}[H]
\centering
\includegraphics[scale=0.6]{../Figures/use-case-diagram-front-end.png}
\caption{Use case diagram - Front-end}
\label{figure:use-case-diagram-front-end}
\end{figure}

\begin{table}[H]
\begin{center}
\begin{tabular}{ l | p{10cm} }
  \hline
  \textbf{Use Case Element} & \textbf{Description} \\ \hline\hline
  Requirements & \hyperref[table:reqfrontend]{FW1}, \hyperref[table:reqip]{FIP3} and \hyperref[table:reqip]{FIP4}\\ \hline
  Application & Front-end \\ \hline
  Name & View data \\ \hline
  Description & Display all the data from the integration platform for the user. \\ \hline
  Actor & User and Integration Platform \\ \hline
  Precondition &
	\par 1. The integration platform has data about the current user.
	\par 2. The integration platform is online.
	\par 3. The user is using a computer with a web browser, that is connected to the internet.
	\\ \hline
  Basic Flow & 
  	\par 1. The user accesses the web page.
  	\par 2. Front-end requests data from the integration platform.
  	\par 3. Integration platform sends data to front-end.
  	\par 4. Front-end displays data to user.
  	\\ \hline
  Alternate Flows & 
  	\par 1A: The user wants to only view heart rate data.
  	\par\hspace{15pt} 1. The user clicks on \textit{Heart Rate} in the navigation bar.
  	\par 1B: The user wants to only view weight data.
  	\par\hspace{15pt} 1. The user clicks on \textit{Weight} in the navigation bar.
  \\ \hline
\end{tabular}
\end{center}
\caption{Textual use case - View data}
\label{table:use-case-view-data}
\end{table}

\subsection{Heart Rate Application}

A use case diagram for the heart rate application is given in figure \ref{figure:use-case-diagram-heart-rate}.
The main objectives of this application is to measure and send the data to the integration platform.

\begin{figure}[H]
\centering
\includegraphics[scale=0.6]{../Figures/use-case-diagram-heart-rate.png}
\caption{Use case diagram - Heart rate application}
\label{figure:use-case-diagram-heart-rate}
\end{figure}

A textual use case for measuring and viewing the heart rate is given in table \ref{table:use-case-measure-heart-rate}.
Figure \ref{figure:use-case-diagram-measure-heart-rate} shows a use case diagram for the textual use case.

\begin{figure}[H]
\centering
\includegraphics[scale=0.75]{../Figures/use-case-diagram-measure-and-view-heart-rate.png}
\caption{Use case diagram - Measure and view heart rate}
\label{figure:use-case-diagram-measure-heart-rate}
\end{figure}

\begin{table}[H]
\begin{center}
\begin{tabular}{ l | p{10cm} }
  \hline
  \textbf{Use Case Element} & \textbf{Description} \\ \hline\hline
  Requirements & \hyperref[table:reqheartrate]{FHR1} and \hyperref[table:reqheartrate]{FHR2} \\ \hline
  Application & Heart Rate Application \\ \hline
  Name & Measure and view heart rate \\ \hline
  Description & Measure the heart rate of the user and then view the measured data. \\ \hline
  Actor & User \\ \hline
  Precondition &
    \par 1. The users phone meets the requirements of the heart rate application.
  	\par 2. The application is installed on the users phone.
  	\par 3. The application is running on the users phone.
  \\ \hline
  Basic Flow & 
  	\par 1. The user places his finger on the camera of the phone.
  	\par 2. The user waits until the heart rate is shown on the display.
  	\par 3. The user views the heart rate.
  \\ \hline
  Alternate Flows & 
  	\par 1A: The user places his finger wrongly, which makes the measured data incorrect.
  	\par\hspace{15pt} 1. The user adjusts his finger.
  \\ \hline
\end{tabular}
\end{center}
\caption{Textual use case - Measure and view heart rate}
\label{table:use-case-measure-heart-rate}
\end{table}

Table \ref{table:use-case-send-heart-rate} gives a textual use case for sending the heart rate to the integration platform.
Figure \ref{figure:use-case-diagram-send-heart-rate} shows a use case diagram for the textual use case.

\begin{figure}[H]
\centering
\includegraphics[scale=0.75]{../Figures/use-case-diagram-send-heart-rate.png}
\caption{Use case diagram - Send heart rate}
\label{figure:use-case-diagram-send-heart-rate}
\end{figure}

\begin{table}[H]
\begin{center}
\begin{tabular}{ l | p{10cm} }
  \hline
  \textbf{Use Case Element} & \textbf{Description} \\ \hline\hline
  Requirements & \hyperref[table:reqheartrate]{FHR3} and \hyperref[table:reqip]{FIP1} \\ \hline
  Application & Heart Rate Application \\ \hline
  Name & Send heart rate \\ \hline
  Description & Send the measured heart rate to the integration platform. \\ \hline
  Actor & User and Integration Platform \\ \hline
  Precondition &
    \par 1. The user has measured his/hers heart rate through the application.
    \par 2. The phone is connected to the internet.
    \par 3. The integration platform is online.
  \\ \hline
  Basic Flow & 
  	\par 1. The user clicks on \textit{Send}.
  	\par 2. The application sends the data to the integration platform.
  	\par 3. The integration platform receives the data.
  \\ \hline
\end{tabular}
\end{center}
\caption{Textual use case - Send heart rate}
\label{table:use-case-send-heart-rate}
\end{table}

\subsection{Weight Application}

The goal of the weight application is to retrieve weight data from HaulthVault and forward it to our integration platform.
Figure \ref{figure:use-case-diagram-weight} shows the main use case diagram for the application.

\begin{figure}[H]
\centering
\includegraphics[scale=0.6]{../Figures/use-case-diagram-weight.png}
\caption{Use case diagram - Weight application}
\label{figure:use-case-diagram-weight}
\end{figure}

In figure \ref{figure:use-case-diagram-view-weight} a use case diagram is given for fetching and viewing the weight data.
Table \ref{table:use-case-view-weight-data} gives a textual use case for the use case diagram.

\begin{figure}[H]
\centering
\includegraphics[scale=0.75]{../Figures/use-case-diagram-view-weight.png}
\caption{Use case diagram - View weight data}
\label{figure:use-case-diagram-view-weight}
\end{figure}

\begin{table}[H]
\begin{center}
\begin{tabular}{ l | p{10cm} }
  \hline
  \textbf{Use Case Element} & \textbf{Description} \\ \hline\hline
  Requirements & \hyperref[table:reqweight]{FHV1} and \hyperref[table:reqweight]{FHV2}\\ \hline
  Application & Weight Application \\ \hline
  Name & View weight data \\ \hline
  Description & Fetch the weight data from HealthVault and display it for the user. \\ \hline
  Actor & User and HealthVault \\ \hline
  Precondition &
    \par 1. The users phone meets the requirements of the weight application.
  	\par 2. The application is installed on the users phone.
  	\par 3. The application is running on the users phone.
  	\par 4. The application is connected to the internet.
  \\ \hline
  Basic Flow & 
  	\par 1. The user writes his email to HealthVault into the email field.
  	\par 2. The user writes his password to HealthVault into the password field.
  	\par 3. The user clicks on a button to login.
  	\par 4. The application shows the data for the user.
  	\par 5. The user views the data.
  \\ \hline
  Alternate Flows & 
  	\par 3A: The user inputs an invalid password or email.
  	\par\hspace{15pt} 1. The user has to go back to step 1 and start over again.
  \\ \hline
\end{tabular}
\end{center}
\caption{Textual use case - View weight data}
\label{table:use-case-view-weight-data}
\end{table}

Table \ref{table:use-case-send-weight-data} gives a textual use case for sending the weight data to our integration platform.
Figure \ref{figure:use-case-diagram-send-weight} displays a use case diagram for the textual use case.

\begin{figure}[H]
\centering
\includegraphics[scale=0.75]{../Figures/use-case-diagram-send-weight.png}
\caption{Use case diagram - Send weight data}
\label{figure:use-case-diagram-send-weight}
\end{figure}

\begin{table}[H]
\begin{center}
\begin{tabular}{ l | p{10cm} }
  \hline
  \textbf{Use Case Element} & \textbf{Description} \\ \hline\hline
  Requirements & \hyperref[table:reqweight]{FHV3} and \hyperref[table:reqip]{FIP2}\\ \hline
  Application & Weight Application \\ \hline
  Name & Send weight data \\ \hline
  Description & Send the weight data from the phone application into the integration platform. \\ \hline
  Actor & User and Integration Platform \\ \hline
  Precondition &
    \par 1. The user has acquired the weight data through the weight application.
    \par 2. The integration platform is online.
  \\ \hline
  Basic Flow & 
  	\par 1. The user presses the \textit{Send} button.
  	\par 2. The Integration platform receives the data.
  \\ \hline
\end{tabular}
\end{center}
\caption{Textual use case - Send weight data}
\label{table:use-case-send-weight-data}
\end{table}
 
% Chapter Template

\chapter{System architecture} % Main chapter title

\label{Architecture} % Change X to a consecutive number; for referencing this chapter elsewhere, use \ref{ChapterX}

\lhead{Chapter 7. \emph{System architecture}} % Change X to a consecutive number; this is for the header on each page - perhaps a shortened title

\section{Overview}

\section{NIPEN}
\section{Front-end}
\section{Heart rate}
\section{Weight}
 
\chapter{Security} 
\label{ch:security}

\lhead{Chapter 6. \emph{Security}}

In this chapter we will discuss the security in our implementation of NIPEN.
We will also talk about what type of security would be expected in a final implementation of NIPEN. 
There will also be a short mention of HIPAA and what role that act will play in a system like this.
At the end of this chapter we will discuss how to grant third party applications limited access to our system.

\section{Security in our solution}
We cleared it with the customer early that to make it easy and achievable for us to make this proof of concept we had to skip most of the security needed.
Therefore the data used is not encrypted. 
We are only using Hypertext Transfer Protocol (HTTP) for transportation of data and not using the Secure Socket Layer (SSL).  
We have also not implemented an authentication system.
Since we have not implemented an authentication system there is only one user of the system.
Every user of the system has access to that users data.

\section{Security in NIPEN}

In this section we will discuss some of the factors that should be considered when implementing NIPEN on a large scale.
We will briefly cover HIPAA and the different forms of security needed for this type of system like authentication, storage and transportation of the data.


\subsection{HIPAA - Health Insurance Portability and Accountability Act}

After discussing security with the customer we brought up the topic of HIPAA \cite{HIPAA}. 
Even though there is no similar act in Norway the customer believed that taking a starting point from HIPAA to discuss the security in NIPEN would be a smart basis. 

\subsection{Storage of data}
When the data is at rest the network should be protected from intrusions.
The hard drives where the data is stored must be encrypted. \cite{Encryption}
HIPAA does not specify what type of encryption should be used, they only specify that it should be encrypted.

\begin{figure}[H]
\centering
\includegraphics[scale=0.50]{../Figures/hipaa.png}
\caption{HIPAA Compliant Data Center Architecture from www.onlinetech.com}
\label{figure:HIPAA}
\end{figure}

\subsection{Transportation of data}
In HIPAA the transportation of data is described that it needs to use Secure Socket Layer. 
This layer is designed to give communication security over the Internet. \cite{SSL}

\subsection{Authentication}
There are multiple systems already in use today in Norway for authenticating users. 
Some of the biggest systems are MinId, BankID, Buypass and Commfides.
The most widespread right now is BankId, with almost 3 million users as of this writing\footnote{\href{https://www.bankid.no/}{https://www.bankid.no/}}.
It would be worth a consideration to use BankID to authenticate the citizens. 
The reason behind this is that it is a widespread system and something the citizens are familiar with. 

When considering what type of authentication the medical professionals should use there are multiple factors to consider. 
Ease of use, efficiency and be able to make an audit trail are some extra factors too consider.
It might be as successful to develop a separate system of authenticating the health professionals if BankID turns out to be cumbersome and unable to be used in the different situations needed for this system. 
For HIPAA there needs to be a system to know which information was accessed, when, and by whom. \cite{Audit}

\subsection{Giving data to third party systems}
Not only should this system accept data from third parties. 
There should also be functionality for exporting data from the system to third party providers. 
For this to happen there needs to be a system for the end users, the citizens, to understand where and to whom the data is going.
Any data that is Protected Health Information \cite{PHI}, which is defined as any data classified as health information connected by the patients identifiers, the end user must control how this information is shared. 
In the end its up to the user who they want to share their data with.
What could be smart is to introduce a Norwegian certification so that the end user knows that the data providers they want to export their data to follows the rules mentioned in previous subsections. 
That way it is easier for the user to trust new systems with their sensitive data.

\section{OAuth 2.0}

OAuth 2.0 is a protocol that could potentially be used for our integration platform.
OAuth 1.0 is outdated, and hence we will mainly focus on OAuth 2.0 in this section.
What this protocol does is that it grants third party applications partial access to a system.
In our case this could for example be only access to the API concerning heart rates or weights in NIPEN.

\subsection{What is OAuth?}

Many applications want to have access to an API provider to gain access to a users information.
However, the user doesn't always want to give his/hers credentials to these applications.
Lets say that a user wants to have an application that is able to upload images to his/hers facebook account.
The user doesn't want to give the application the credentials for facebook.
This would be highly insecure, since once the application has the authentication information of a user it is able to do whatever it wants with the account.
To handle this we need some kind of protocol, this is where OAuth comes in.
With OAuth we are able to give a third party application partial access to a system without handing over the credentials of a user.
This is a token based system, where the application receives a token which grants it access to use some of the systems services.

When a third party app wants to use a system that uses OAuth, it first needs to be registered in that system.
Under registration the system usually wants some basic information about the application.
This might include name of the application, a description, logo and etc.
A redirect URI also needs to be registered.
This URI needs to use TLS, i.e. must begin with \textit{https}.
This address is used to send an authentication token from the system to the application.
After the registration, the application should receive a client ID which is used to identify the application at the system.

Now users should be able to connect the application with theirs account on the system.
The way it works is that through the application they will be sent to a website of the system the application wants to connect with.
There the user will be asked if he/she wants to grant permissions to the application.
If the user agrees, then he/she will be redirected to the applications specified redirect URI with an access token as a parameter.
The application will now use this token to access the system.
When the user doesn't want the application to have access to the system any more, he/she can simply disallow the applications access through the systems website.

An example of how OAuth works is given in figure \ref{figure:oauth-in-a-nutshell}.
In this example we see a website that wants to send images to a users facebook account.
If the user allows the application access through facebook, then facebook will send an access token to the website.
Then \textit{www.ImageToFacebook.com} is able to send images to the users account, where the access token is used for authentication.
The user should be able to disallow the website access at any time on facebooks website. \cite{OAuth}

\begin{figure}[h]
\centering
\includegraphics[scale=1.0]{../Figures/oauth-in-a-nutshell.png}
\caption{OAuth in a nutshell}
\label{figure:oauth-in-a-nutshell}
\end{figure}

\iffalse
\subsection{Applying OAuth to NIPEN}

It should not be a problem to apply OAuth to NIPEN.
The functionality that NIPEN needs to implement is a capability of registering third party applications.
Whenever third party applications are registered, they should be associated with a client ID.
This ID should be sent to the third party application.
Users should be able to allow different third party applications limited access to their account.
This could for example be access to only send heart rate measurements (or weights).
When a user allows access to an application an access token should be sent to the specified application, so it gets partial access to the users NIPEN account.
At any time users should be able to remove the applications access to NIPEN.

A summary of how this would work is given in the example below, which illustrates how the granting of permission would work with our heart rate application:

\begin{enumerate}
\item User starts the heart rate application.
\item User wants to grant the application access to push heart rate measurements to NIPEN.
\item User is redirected to the website of NIPEN.
\item User logs in to NIPEN (with e.g. BankID).
\item NIPEN asks the user if he wants to grant the heart rate application access to push heart rate data into NIPEN.
\item If user agrees, an access token is sent to the heart rate application.
\end{enumerate}

With OAuth we should also expand the JSON messages that are sent to NIPEN.
They should now contain a client ID and an access token.
The client ID is used to identify the application that is pushing values to NIPEN, while the access token is used to authorise the application.
A JSON string for a heart rate measurement could now look something like this:

\begin{lstlisting}[language=json]
{
	"clientID": 489431,
	"accessToken": "safDFSAadsffsasdFDSewfaDSFAdsfaewRETrehhgreeErw",
	"userId":453,
	"timestamp":"2013-11-05 12:12:38.0",
	"value":66,
	"unit":"bpm"
}
\end{lstlisting}

This information should be sent in an encrypted format by using https.

NIPEN will now be able to identify what application is sending data, and is also capable of checking if the application is authorised.
The user is also able to grant other applications partial access to his/hers account, without giving these applications their credentials.
Hence, this should be a good way of handling third party applications in NIPEN.

\fi 

\chapter{Sprint 0}
\label{Sprint0}
\lhead{Chapter 6. \emph{Sprint 0}}

\section{Duration}
The duration of the sprint was the following:
\begin{itemize}
\item Sprint start: August, 26th
\item Sprint end: September, 8th
\end{itemize}

\section{Planning}

Our first sprint was focused on project planning and preliminary studies.
We planned to receive some kind of approval on our understanding of the project
as described by the customer before starting preliminary studies. This was done mainly to avoid
\begin{enumerate}[a)]
\item planning a solution for a different problem and
\item wasting resources studying something that would not be relevant for the project.
\end{enumerate}
In an effort to prepare a draft of system requirements to be approved by the customer,
we planned some brainstorming sessions to pin down a number candidates for health measurements
(weight, blood pressure\ldots) to be supported by the system.

\section{Goal(s)}

We wanted to achieve a sufficient understanding of the problem's scope,
relevant technologies and similar solutions as well as:% complete a high-level planning which would have included:
\begin{enumerate}[a)]
\item identify high-level requirements, phases and milestones
\item identify a high-level system architecture
\item choose a development process
\end{enumerate}

Additionally, we wanted to decide what software to use to write the report,
choose and setup development enviroments, SDKs and version control systems in order to be able
to get going with development and documentation as soon as necessary.

\begin{figure}[h]
\centering
\includegraphics[scale=0.60]{../Figures/burndownSprint0.png}
\caption{Burndown chart Sprint 0}
\label{figure:burndownsprint0}
\end{figure}

\newpage
\section{Backlog}

See below the sprint backlog.
\begin{enumerate}[1.]
	\item \textbf{Project planning} for this sprint included:
		\begin{itemize}
			\item \textbf{Booking of rooms}:
				where to hold meetings and work.
			\item \textbf{Scheduling weekly meetings}:
				with both the customer and supervisor.
			\item \textbf{Choice of development process}:
				candidates are Waterfall and Scrum.
			\item \textbf{Requirements gathering}:
				an initial set (draft) of high-level requirements.
			\item \textbf{Risk analysis}:
				a first draft.
		\end{itemize}
	\item \textbf{Project management} included:
		\begin{itemize}
			\item \textbf{Sprint startup meeting}:
				included sprint planning and review
			\item \textbf{Weekly meetings}:
				with both customer and supervisor (from the 26th).
			\item \textbf{Meeting notes}:
				taking notes during meetings, reviewing of the notes.
			\item \textbf{Status reports}:
				for weeks 35 and 36.
		\end{itemize}
	\item \textbf{Read course's compendium}
	\item \textbf{Preliminary studies}\newline
		of relevant technologies like databases, web-servers and web APIs.
	\item \textbf{Identify possibile measurements}\newline
		to be supported by the system (glucose, weight\ldots).
	\item \textbf{Preliminary system architecture}\newline
		brainstorm on a high-level architecture to be submitted to the client for feedback.
	\item \textbf{Choice and setup of project's tools}\newline
		such as tools for the report and version control system.
	\item \textbf{Android environment setup}\newline
		both IDEs and SDK.
	\item \textbf{Group dynamics lecture}\newline
		attendance is mandatory.
\end{enumerate}


\section{Results and feedback}

\begin{figure}[h]
\centering
\includegraphics[scale=0.30]{../Figures/architecture-draft.png}
\caption{System architecture (draft)}
\label{figure:architecture-draft}
\end{figure}

During the first days of the project, which were technically not part of any sprint,
we met each other, the customer and the supervisor.
The customer described the problem, explaining what he had in mind and why it was relevant for him.
The supervisor handed us some notes to help the project kick-off.
We exchanged contacts information (both mail and Skype) to make sure everybody could reach out to
everybody else. Additionally, we set up a Facebook group where to post relevant informations
and a Google Drive folder to share and collaborate on editing documents.
We scheduled weekly meetings with all involved parties as follows:
\begin{itemize}
\item customer and supervisor meetings on Mondays
\item collective work sessions on Mondays (10.00 - 19.00) and Wednesdays (15.00 - 19.00)
\end{itemize}

At the beginning of the sprint we shared our thoughts on what we had understood from the customer.
In order to ensure we had a proper understanding of the problem we produced a diagram
(figure \ref{figure:architecture-draft}) representing a high-level architecture
of the product (an \textit{ehealth integration platform}) and submitted it to the customer
via mail for feedback.

It consisted of:
\begin{enumerate}[1.]
\end{enumerate}

We then begun studying relevant technologies and similar solutions and eventually identified
a number of candidate health measurements and submitted those to the customer as well.

Our idea was to develop three prototype applications, other than the integration
platform and front-end, to be able to demonstrate how these systems interoperate.
Each of these prototypes would have focused on a single health measurement.

We found an agreement with the customer on heart rate and weight measurements which we had
proposed based on two facts. Firstly, we had found an open-source Android application which could
measure the user's heart rate using the mobile's camera and hoped we could use it for one prototype.
Secondly, one team member owned a Withings weighing scale device which we tought
would be interesting to use for another prototype.
The customer agreed on these candidates because they were two of the most common measurements.

A possible third application was discussed both internally and with the customer but
given time and resource constraints of this project we intentionally left it rather vague
as we were unsure whether we would have had time for that.

By the end of the sprint the customer started mentioning Microsoft's HealthVault and suggesting
that some kind of interoperability with HealthVault would have proved valuable for him.

With regards of our development process, we decided to use Scrum based on various considerations as
described in \ref{section:development-methodology}. To better perform Scrum, we decided
to use an online tool called \textit{Scrumdo}, described in section \ref{section:tools}.
We also chose and setup other tools like:
\begin{itemize}
\item version control system (Git and GitHub)
\item development environment (IntelliJ IDEA)
\item documentation (LaTeX)
\end{itemize}

\section{Evaluation}

%The sprint generally achieved its expected results.
At the end of the sprint we had achieved a good understanding of the problem and a good overview on a
high-level system architecture and requirements. We didn't get as far as we hoped on our studies on some particular
technologies and this was due to a number of factors. Firstly, estimating the amount of time required to obtain
a sufficient knowledge in a topic resulted difficult and we generally needed more time than we thought and had planned.
Secondly, we had little familiarity with most involved technologies.
Given the size of our group, we had some difficulties mitigating this problem due to
our reduced technical skill pool. Some difficult topics of study were Apache Camel, its extension Apache IPF,
and whether it was desiderable to use these instead of the Spring framework (and why).
Although by the end of the sprint we came to an agreement on what technologies we need to use
(database, web server\ldots) we didn't decide on a specific implementation of such technologies
(Apache Tomcat, Spring, Apache Camel, JAX-RS\ldots).

During this sprint we did not manage to send all the documents the supervisor
required or they were missing some important detail. However, what we sent
was sufficient to give him a good understanding of our progress so we
consider this a minor issue.
 


\chapter{Sprint 1}
\label{Sprint0}
\lhead{Chapter 7. \emph{Sprint 1}}

\section{Goal(s)}
This sprint's goal was the completition of an initial prototype of the system whose functionality
had to be demonstrated to the customer; this corresponded to project milestone M1.
The purpose of such prototype was to act as a proof-of-concept and give the customer
a chance to express some feedback upon which we hoped to start a discussion
on possible improvements and new features.

\section{Planning}
In order to achieve the goal for this sprint we planned to accomplish an initial design
and implementation of every part of the system including:
a) the frontend b) the backend and database c) the Android application to measure heart rate.
We planned some additional studies on relevant technologies for the first week of the sprint
and begin writing some parts of the report.

\section{Duration}
The duration of the sprint was the following:
\begin{itemize}
\item Start: September, 9th
\item Milestone M1 (first system prototype): September, 20th
\item End: September, 22nd
\end{itemize}

\begin{figure}[H]
\centering
\includegraphics[scale=0.60]{../Figures/burndownSprint1.png}
\caption{Burndown chart Sprint 0}
\label{figure:burndownsprint0}
\end{figure}

\section{Backlog}

See below the sprint backlog.
\begin{itemize}
	\item \textbf{M1 First system prototype}
	\item \textbf{Project management}\newline
	This included:
	\begin{itemize}
		\item \textbf{Weekly startup meeting}: 
		\item \textbf{Meeting notes}:
			taking notes during meetings, reviewing of the notes.
		\item \textbf{Status reports}:
			for both week 37 and 38
		\item \textbf{Risk analysis}:
			updated on a weekly basis, so twice per sprint.
			The risk analisys was submitted to the supervisor and the customer.
		\item \textbf{Planning for the next iteration}:
			the project manager prepared a plan for the next iteration
			which would be illustrated and agreed upon on next iteration's startup meeting.
	\end{itemize}
	\item \textbf{Weekly meetings}
		includes meetings with both the customer and the supervisor.
		The meeting with the customer was held on Skype.
	\item \textbf{Additional pre-studies}
		Continued studies on relevant technlogies such as:
	\begin{itemize}
		\item \textbf{HealthVault}: Microsoft's online health platform.
		\item \textbf{Apache Camel}: a routing engine for enterprise integration patterns.
		\item \textbf{Javascript libraries for charts}: to be used in the frontend.
	\end{itemize}
	\item \textbf{System development}
		Initial design and implementation. For this sprint, this accounted for:
	\begin{itemize}
		\item \textbf{Backend development}:
			Development of Spring controllers for API endpoints. Database (DAO) coding to
			implement data persistence.
		\item \textbf{Frontend development}:
			Coding of the frontend. This included setting up an HTML page which used
			AJAX to perform API calls and JS to show the data retrieved using a bar chart.
		\item \textbf{Deployment}:
			Deployment of both backend and frontend using a servlet container (Tomcat).
	\end{itemize}
	\item \textbf{Heart rate application}:
		basic implementation of the Heart rate application. The application should be able to acquire
		the user's heart reate and send perform an API call to store the data on the backend.
	\item \textbf{Database development}:
		We chose a suitable database to use for implementing persistency on the backend.
		We opted for MySQL due to the familiarity we had with it. We then deployed it on the server machine
		and added a first table so that we could test the functionality of the backend.
	\item \textbf{Testing}:
		Perform unit and integration testing for the heart rate application and the backend.
\end{itemize}


%% i actually dont think having a table for this is a good idea. tables can't have much information
%% also it would require a lot of tweaking to get the estimated/actual times to look reasonable.
%% maybe a textual description would be better so we can omit some details in favor of others.
\iffalse
\begin{table}
\begin{tabular}{ | l | l | l | l | }
 \hline
  Story ID & Description & Size & Assignee \\
  \hline\noalign{\smallskip}\noalign{\smallskip}\hline
  33 & Project Management			& 8	& Emanuele  \\
  12 & M1 First System prototype	& 0 & All		\\
  45 & Weekly meetings (week 37)	& 6 & All		\\
  42 & Additional pre-studies		& 5 & Emanuele	\\
  \hline
\end{tabular}
\caption{}
\label{}
\end{table}
\fi

\section{Testing}

\section{Results and feedback}

Around the end of this sprint, on September 20th, we demonstrated a prototype of system to
the client which included: a) the frontend, b) the backend, c) an Android application to measure
heart rate using the phone's camera.

The backend supported web API calls for storing and retrieving heart rate measurements.
Both the Android application and the frontend used these API calls for sending and retrieving
such measurements respectively. The measurements were ultimately shown to the user on a
webpage (the frontend) using a bar chart.

The customer was pleased with the results, stating that they were above his expectations.

We discussed about which features he would like to see implemented next in the product
and we agreed to prioritise interoperability with HealthVault instead of Withings.

The reason behind this decision was the fact that interoperability with HealthVault would have
enabled the product to use third-party devices supported by HealthVault to gather health data.

Since the number of devices supported by HealthVault is substantial, this was deemed a
desiderable feature for the product.

The customer asked if the amounts of resources we had at disposal to actually implement
such interoperability was sufficient and inquired about the general viability of such approach.

We expressed our confidence about the feasibility of such approach and that the resources
at hand were sufficient. Nevertheless, we set a deadline (10 days) to assess such
feasibility in order not to delay the general progress of the product.

\textbf{Notes}: before the end of the sprint one team member had moved permanently to Oslo.

we didn't manage to begin writing the report.

\section{Evaluation}

The sprint was successful as we managed to achieve the main goal of the sprint (Milestone M1).
We were pleased by the positive feedback received from the customer and felt motivated to keep up the good work.
Furthermore, having proactively partecipated in the improvement/proposal of product requirements
together with the customer made the whole team look forward to implement these new features and improve the product.

Although we didn't manage write anything in the report we still had
 


\chapter{Sprint 2}
\label{Sprint0}
\lhead{Chapter 8. \emph{Sprint 2}}

\section{Goal(s)}
Having just demonstrated the product to the customer in the previous sprint the main goal
for this iteration was to write the report in view of the deadline for mid-term delivery.
Additionally we had to assess the feasibility of the interoperability of our system with HealthVault.


\section{Planning}
Most of the work planned for this sprint was documentation writing.
This included writing the report itself but also reviewing other documents such as the templates, notes and agendas.

This plan made it easy to assign working tasks despite the fact that one member
of the team was working on the project remotely from another city.

Report work for this sprint was in fact something that we could


We also planned to begin some preliminary work toward the second prototype of the system which would
feature HealthVault interoperability. Our plan was to develop an Android application which would connect
to HealtVault in order to retrieve data about user's weight and forward it to our system. For this task
we planned to use HealthVault's SDK for Android.


\section{Duration}
\begin{itemize}
\item Start: September, 23rd
\item End: October, 6th
\end{itemize}

\section{Backlog}

\begin{itemize}
	\item \textbf{Project management}\newline
	This included:
	\begin{itemize}
		\item \textbf{Weekly startup meeting}:
		\item \textbf{Meeting notes}:
			taking notes during meetings, reviewing of the notes.
		\item \textbf{Status reports}:
			for both week 39 and 40
		\item \textbf{Review templates}:
			we produced some templates to be adopted for status reports and other documents.
		\item \textbf{Risk analysis}:
			updated on a weekly basis, so twice per sprint.
			The risk analisys was submitted to the supervisor and the customer.
		\item \textbf{Planning for the next iteration}:
			the project manager prepared a plan for the next iteration
			which would be illustrated and agreed upon on next iteration's startup meeting.
	\end{itemize}
	\item \textbf{Weekly meetings}
		included meetings with both the customer and the supervisor.
		*During week 39.we were unable to contact the customer. This was not a problem*
	\item \textbf{Report work}:
		we started writing the report.
		This tasks was definitely the biggest of the sprint and accounted for the majority of the work.
		We focused on writing the abstract and first chapters required for the mid-term delivery.
	\item \textbf{HealthVault studies}:
		we performed further studies in order to assess the feasibility of developing a prototype
		application that implemented interoperability with Microsoft's HealthVault platform.
	\item \textbf{System development}:
		this included some preliminary work on the backend to accomodate for the new weight API.
	\item \textbf{Application development}:
		Initial design and implementation of an Android application using the weight API.

\end{itemize}



Make an overall project plan (Gantt diag.)	emanuele	Done

Review system requirements	emanuele	Done
Review templates	Done


\section{Results and feedback}

Results for this sprint are positive. We managed to write the initial chapters of the report.
Also, studies on HealtVault went smoothly and we had a good start at developing the Android application
especially. In fact, using HealthVault's SDK greatly reduced the amount of code we needed to write ourselves
from scratch.
System development also proceeded smoothly and we added support

\section{Evaluation}

our thoughts about this sprint 

\chapter{Sprint 3}
\label{Sprint3}
\lhead{Chapter 10. \emph{Sprint 3}}

\section{Duration}
The duration of the sprint was the following:
\begin{itemize}
\item Sprint start:  October, 7th
\item Milestone M2 (mid-term report delivery): October, 14th.
\item Sprint end: October, 20th
\end{itemize}

\section{Planning}

The last sprint had proceeded without issues, so we didn't make any big adjustments to our plan for this iteration. 
We planned to focus on report work and do some more development after achieving milestone M2.

We understood that mid-term delivery was a good chance to receive valuable feedback from the student advisor, the technical advisor and the customer so we directed most of our efforts in report writing in order to receive as much feedback as possible on it.

We prepared a separate work plan for our colleague in Oslo with a focus on documentation
writing since we thought it would be easier for him and for us to work collaboratively
on the report rather than development.

\iffalse
The first week the whole team would focus on writing the report
%The plan consisted on report writing for the first week,
in order to receive as much feedback as possibile on it. 
%This task was assigned to the whole team.
For the second week, %in view of our customer's visit to Trondheim on the last days of the month,
we planned to proceed with system and application development %after mid-term delivery
in order to implement most of the functionality for the second product prototype meanwhile
our colleague in Oslo would continue to write the report.
%as well as some additional work on the report to be done by our colleague in Oslo.

%assigning him some work on the report for the second week of the sprint as well.
\fi

\section{Goal(s)}

The goal of this sprint was the mid-term delivery of the report; this corresponded to
the second project's milestone M2. Furthermore, another goal was finalizing the design and continuing
the development of the second prototype of the product.% which would feature HealthVault interoperability.


\section{Results and feedback}

We were able to write what we planned to in the report and had no problem meeting the deadline for the midterm delivery. 
We were generally satisfied with the quality of what had been written having spent some time reviewing the document to ensure consistency and a logical flow.
We attended the lecture about technical writing and at the end of which we received some positive feedback on the report.

Development proceeded without major issues. 
We were able to re-use good portions of code provided by HealthVault's SDK, dramatically reducing the time spent on development.
We performed some major refactoring of the code on a dedicated Git branch in an effort to a achieve a more \'neat\' deployment solution which didn't require a separate servlet container like Apache Tomcat but was entirely based on Spring, using an embedded servlet-container.

During one of the meetings with the customer, we scheduled a demonstration of the product for the 24th of October when he would travel to Trondheim for work.

\section{Evaluation}

\begin{figure}
\centering
\includegraphics[scale=0.60]{../Figures/burndownSprint3.png}
\caption{Iteration burndown chart}
\label{figure:burndownsprint3}
\end{figure}

All in all this sprint was successful because we managed to meet the deadline for report delivery.
Additionally, we were glad to have received positive feedback on our report.

However, refactoring was an unexpectedly time-consuming task and did not achieve expected results.
Luckily, since all refactoring had been performed on a separte branch on Git, it didn't create problems as we always had a working branch for development and demonstration purposes.

By the end of the sprint it became clear that the customer wasn't that interested in a prototype application for Withings. 
However, because one team member owned a Withings scale device, we thought it would have been a good idea to include that in the demonstration.
Since Withings supported integration with HealthVault natively, we came up with the following scenario for product demonstration:
\begin{enumerate}[1)]
\item the user steps on the scale
\item the weight is acquired by the scale and sent to Withings
\item the measurement on Withings is automatically integrated in HealthVault
\item the Android application fetches the measurement from HealthVault
\item the measurement is sent to our integration platform
\end{enumerate}
Regarding our collaboration with our colleague in Oslo, he didn't report the time he spent on the project correctly, and his contribution was discontinuos because he had begun to work full-time.
This didn't create any big problems because development had proceeded smoothly so other team members had some spare time to cover up for that.
This iteration burndown chart is shown in figure \ref{figure:burndownsprint3}.

\clearpage
\section{Backlog}
See below the sprint backlog.

\begin{enumerate}[1.]
\item \textbf{M2 Mid-term report}:
	delivery of the mid-term report
\item \textbf{Project management} included:
	\begin{itemize}
		\item \textbf{Sprint startup meeting}:
			included sprint planning and review
		\item \textbf{Weekly meetings}: 
			with both the customer and the supervisor. Internal meeting with our colleague in Oslo
		\item \textbf{Meeting notes}:
			taking notes during meetings, reviewing of the notes
		\item \textbf{Status reports}:
			for both week 41 and 42
		\item \textbf{Risk analysis}:
			updated on a weekly basis
	\end{itemize}
	\item \textbf{Weighter Application development}\newline
		continued development of the Android application
	\item \textbf{System development}\newline
		study different ways to deploy the backend
		1) as a WAR file in a separate servlet container (Apache Tomcat) or 2) as a Spring JAR file
		with embedded servlet support. Furthermore we continued development
	\item \textbf{Report work}\newline
		to be done by the team member in Trondheim
	\item \textbf{Additional report work}\newline
		to be done by the team member in Oslo
	\item \textbf{Testing}\newline
		perform unit testing
\end{enumerate}
 


\chapter{Sprint 4}
\label{Sprint0}
\lhead{Chapter 10. \emph{Sprint 4}}

\section{Goal(s)}
what did we expect to achieve by the end of this sprint (general progress in the project)

\section{Planning}
what we planned to do
\section{Duration}
\section{Backlog}
include sprint backlog
\section{Results and feedback}
from customer, from supervisor
\section{Evaluation}
our thoughts about this sprint 


\chapter{Sprint 5}
\label{Sprint0}
\lhead{Chapter 11. \emph{Sprint 5}}

\section{Goal(s)}

This was the final sprint of our project. Its main goal was to
finalize the product and the documentation as well as prepare
a presentation.

\section{Duration}
The duration of the sprint was the following:
\begin{itemize}
\item Sprint start:  November, 4th
\item Milestone M3 (feature freeze): November, 4th.
\item Sprint end: November, 20th
\end{itemize}

\begin{figure}[h]
\centering
\includegraphics[scale=0.60]{../Figures/burndownSprint5.png}
\caption{Burndown chart Sprint 5}
\label{figure:burndownsprint5}
\end{figure}

This sprint was a bit longer than the others because it included the last days
that were left till the end project.

\section{Planning}

The plan for this sprint consisted primarily in report writing and review.
At the beginning of the sprint we agreed with the supervisor to deliver him a mostly
finished version of the report by the end of the first week in order to let him
give us some last feedback on it.

\section{Backlog}
See below the sprint backlog.
\begin{enumerate}[1.]
\item a
\end{enumerate}

\section{Results and feedback}

from customer, from supervisor
\section{Evaluation}
our thoughts about this sprint 

\chapter{Testing}

\label{Testing}

\lhead{Chapter 12. \emph{Testing}}

\section{Testing methods}

\section{Tools}

\section{Functional tests}

\begin{table}
\begin{center}
\begin{tabular}{ | l | p{10cm} | }
	\hline
	\textbf{Test}	&	\textbf{ID 1} \\
	\hline\noalign{\smallskip}\noalign{\smallskip}\hline
	Name				& Heart rate REST (in) \\
	Requirement			& FIP1 \\
	Description			& Test the REST endpoint for receiving heart rate data models \\
	Preconditions		& 	\par The IP is deployed and running on a server machine 
							\par The server machine is accesible from the machine running the test \\
	Steps 				&	\par 1. Enter username
							\par 2. Enter password
							\par 3. Press login button \\
	Postconditions		& The table \textit{heart rate} in the database on the server machine contains a row \\
	Results				& - \\
	Comments			& - \\
	Status				& OK or FAIL \\
	Tester				& Person \\
	Date				& dd-mm-2013 \\
	\hline
\end{tabular}
\end{center}
\end{table}

\begin{table}
\begin{center}
\begin{tabular}{ | l | p{10cm} | }
	\hline
	\textbf{Test}	&	\textbf{ID 2} \\
	\hline\noalign{\smallskip}\noalign{\smallskip}\hline
	Name				& Weight REST (in) \\
	Requirement			& FIP2 \\
	Description			& Test the REST endpoint for receiving weight data models \\
	Preconditions		&	\par The IP is deployed and running on a server machine
							\par The server machine is accesible from the machine running the test \\
	Steps 				&	\par 1. Enter username
							\par 2. Enter password
							\par 3. Press login button \\
	Postconditions		& The table \textit{weight} in the database on the server machine contains a row \\
	Results				& - \\
	Comments			& - \\
	Status				& OK or FAIL \\
	Tester				& Person \\
	Date				& dd-mm-2013 \\
	\hline
\end{tabular}
\end{center}
\end{table}

\begin{table}
\begin{center}
\begin{tabular}{ | l | p{10cm} | }
	\hline
	\textbf{Test}	&	\textbf{ID 3} \\
	\hline\noalign{\smallskip}\noalign{\smallskip}\hline
	Name				& Heart rate REST (out) \\
	Requirement			& FIP3 \\
	Description			& Test the REST endpoint for requesting heart rate data models \\
	Preconditions		&	\par The IP is deployed and running on a server machine
							\par The server machine is accesible from the machine running the test
							\par At least one heart rate measurement has been stored on the IP \\
	Steps 				&	\par 1. Enter username
							\par 2. Enter password
							\par 3. Press login button \\
	Postconditions		& The login screen disappear and the inbox appear \\
	Results				& Correct \\
	Comments			& - \\
	Status				& OK \\
	Tester				& Person \\
	Date				& 21-10-2013 \\
	\hline
\end{tabular}
\end{center}
\end{table}


\begin{table}
\begin{center}
\begin{tabular}{ | l | p{10cm} | }
	\hline
	\textbf{Test}	&	\textbf{ID 4} \\
	\hline\noalign{\smallskip}\noalign{\smallskip}\hline
	Name				& Weight REST (out) \\
	Requirement			& FIP4 \\
	Description			& Test the REST endpoint for requeting weight data models \\
	Preconditions		&	\par The IP is deployed and running on a server machine
							\par The server machine is accesible from the machine running the test
							\par At least one weight measurement has been stored on the IP \\
	Steps 				&	\par 1. Enter username
							\par 2. Enter password
							\par 3. Press login button \\
	Postconditions		& The login screen disappear and the inbox appear \\
	Results				& Correct \\
	Comments			& - \\
	Status				& OK \\
	Tester				& Person \\
	Date				& dd-mm-2013 \\
	\hline
\end{tabular}
\end{center}
\end{table}


\begin{table}
\begin{center}
\begin{tabular}{ | l | p{10cm} | }
	\hline
	\textbf{Test}	&	\textbf{ID 5} \\
	\hline\noalign{\smallskip}\noalign{\smallskip}\hline
	Name				& Heart rate measurement \\
	Requirement			& FHR1 \\
	Description			& Test the heart rate measurement functionality \\
	Preconditions		& The application has started \\
	Steps 				&	\par 1. User holds his finger on the camera applying a slight pressure \\
	Postconditions		&	\par 1. The icon on the left side of the screen is blinking 
							\par 2. A measurement appears after 4 seconds at most \\
	Results				& Correct \\
	Comments			&	The measurement is expected to be rough.
							Any value between 60 and 100 is okay as long as it actually varies slightly based
							on the tester's perceived heart rate.  \\
	Status				& OK \\
	Tester				& Person \\
	Date				& 21-10-2013 \\
	\hline
\end{tabular}
\end{center}
\end{table}


\begin{table}
\begin{center}
\begin{tabular}{ | l | p{10cm} | }
	\hline
	\textbf{Test}	&	\textbf{ID 6} \\
	\hline\noalign{\smallskip}\noalign{\smallskip}\hline
	Name				& Login \\
	Requirement			& FR1 \\
	Description			& Test the login functionality \\
	Preconditions		& The application has started \\
	Steps 				&	\par 1. Enter username
							\par 2. Enter password
							\par 3. Press login button \\
	Postconditions		& The login screen disappear and the inbox appear \\
	Results				& Correct \\
	Comments			& - \\
	Status				& OK \\
	Tester				& Person \\
	Date				& 21-10-2013 \\
	\hline
\end{tabular}
\end{center}
\end{table}


\begin{table}
\begin{center}
\begin{tabular}{ | l | p{10cm} | }
	\hline
	\textbf{Test}	&	\textbf{ID 7} \\
	\hline\noalign{\smallskip}\noalign{\smallskip}\hline
	Name				& Heart rate measurement \\
	Requirement			& FHR3 \\
	Description			& Test the heart rate application interoperability \\
	Preconditions		&	\par 1. The application has started
							\par 2. A heart rate measurement has been acquired \\
	Steps 				&	\par 1. User presses \textbf{Send} button \\
	Postconditions		&	\par 1. The application shows a toast
							\par 2. The measurement is acquired by the Integration Platform\\
	Results				& Correct \\
	Comments			& --\\
	Status				& OK \\
	Tester				& Person \\
	Date				& 21-10-2013 \\
	\hline
\end{tabular}
\end{center}
\end{table}

\begin{table}
\begin{center}
\begin{tabular}{ | l | p{10cm} | }
	\hline
	\textbf{Test}	&	\textbf{ID 8} \\
	\hline\noalign{\smallskip}\noalign{\smallskip}\hline
	Name				& HealthVault data fetching \\
	Requirement			& FHV1 \\
	Description			& Test HealthVault connectivity \\
	Preconditions		&	\par 1. The application has started
							\par 2. The user has authenticated to HealthVault
	Steps 				&	\par 1. User \\
	Postconditions		&	\par 1. The icon on the left side of the screen is blinking.
							\par 2. A measurement appears after 4 seconds at most.\\
	Results				& Correct \\
	Comments			&	The measurement is expected to be rough.
							Any value between 60 and 100 is okay as long as it actually varies slightly based
							on the tester's perceived heart rate.  \\
	Status				& OK \\
	Tester				& Person \\
	Date				& 21-10-2013 \\
	\hline
\end{tabular}
\end{center}
\end{table}

\begin{table}
\begin{center}
\begin{tabular}{ | l | p{10cm} | }
	\hline
	\textbf{Test}	&	\textbf{ID 5} \\
	\hline\noalign{\smallskip}\noalign{\smallskip}\hline
	Name				& Heart rate measurement \\
	Requirement			& FHR1 \\
	Description			& Test the heart rate measurement functionality \\
	Preconditions		& The application has started \\
	Steps 				&	\par 1. User holds his finger on the camera applying a slight pressure. \\
	Postconditions		&	\par 1. The icon on the left side of the screen is blinking.
							\par 2. A measurement appears after 4 seconds at most.\\
	Results				& Correct \\
	Comments			&	The measurement is expected to be rough.
							Any value between 60 and 100 is okay as long as it actually varies slightly based
							on the tester's perceived heart rate.  \\
	Status				& OK \\
	Tester				& Person \\
	Date				& 21-10-2013 \\
	\hline
\end{tabular}
\end{center}
\end{table} 
\chapter{Conclusion}
\label{ch:conclusion}
\lhead{Chapter 14. \emph{Conclusion}} 

In this chapter we provide a conclusion for this work and outline some
starting points for further work.

\section{Conclusion}

We were definitely happy with the project's task which led us to study trending topics
like eHealth and self-quantification. Although we had no previous experiences in designing and developing a web API,
we found it interesting and enjoyable. This resulted in a proactive participation by the group
in the improvement and proposal of requirements, a process which in turn made the project even more interesting
for us to work on.

Preliminary studies have shown a number of existing solutions



Our architectural design choices for the integration platform make it well suited to be scaled up.
The integration platform in fact, consists of a web application and a database which can be deployed separately,
reducing the amount of resources needed on a single machine.
The same applies to the front-end as well because it interacts with the integration platform using its API just like
any other application and therefore does not rely on being deployed on the same machine.
Furthermore, because we are not employing a classical SOAP client-server architecture but a more lean, REST-like one,
the implementation can be considered lightweight in term of resources and it is easier for application developers to interact with.

\section{Further work}

There are certainly many limitation to our product, as described in \ref{subsec:limitations}, which we believe should be
the focus of further work. In this section we present our suggestions on how to overcome said limitations.

\textbf{API design}
A functional system should have a much more extensive and complete API which should be
thoroughly documented and accompanied by a number of examples and libraries which
would contribute to reduce the development effort for third-parties.
Data models should be complete and able to represent a meaningful amount of information,
and possibly exhibit some degree of interoperability with other established standards
such as HL7\footnote{www.hl7.org}, DICOM and CEN's TC/251.

\textbf{Deployment}
The actual architecture makes it possible with some slight refactoring to simply separate
the codebase of the backend and the frontend and deploy them individually on separate machines.
Also, connecting to a database on a different machine can be achieved by simply
modifying the connection address.

\iffalse \textbf{Authentication} \fi


\textbf{Authorization}
Our suggested approach is to implement an OAuth\footnote{www.oauth.net} based authorization mechanism
for third-parties applications. Users should be able to allow different third party applications limited access
to their account, for example only to weights measurements.
%It should not be a problem to apply OAuth to NIPEN.
%The functionality that NIPEN needs to implement is a capability of registering third party applications.
Whenever third party applications are registered, they should be associated on the system with a client ID which
should then be returned to the the application itself.
When a user allows access to an application an access token representing the level of authorization should be sent
to the specified application, so it gets partial access to the users NIPEN account.
At any time users should be able to remove the applications access to NIPEN.

A summary of how this would work is given in the example below, which illustrates how \iffalse the granting of
permission\fi authorization would work with our heart rate application:

\begin{enumerate}
\item User starts the heart rate application.
\item User wants to grant the application access to push heart rate measurements to NIPEN.
\item User is redirected to the website of NIPEN.
\item User logs in to NIPEN (with e.g. BankID).
\item NIPEN asks the user if he wants to grant the heart rate application access to push heart rate data into NIPEN.
\item If user agrees, an access token is sent to the heart rate application.
\end{enumerate}

With OAuth we should also expand the JSON messages that are sent to NIPEN.
They should now contain a client ID and an access token. The client ID is used to identify the application
\iffalse that is pushing values to NIPEN\fi, while the access token
is used for authorization. A JSON string for a heart rate measurement would now look like this:

\begin{lstlisting}[language=json]
{
	"clientID": 489431,
	"accessToken": "safDFSAadsffsasdFDSewfaDSFAdsfaewRETrehhgreeErw",
	"userId":453,
	"timestamp":"2013-11-05 12:12:38.0",
	"value":66,
	"unit":"bpm"
}
\end{lstlisting}

This information should be sent in an encrypted format by using https.

NIPEN will now be able to identify what application is sending data, and is also capable of checking if the application is authorized.
The user is also able to grant other applications partial access to his/hers account, without giving these applications their credentials.
Hence, this should be a good way of handling third party applications in NIPEN.
 
\chapter{Evaluation} 
\label{ch:evaluation}

\lhead{Chapter 13. \emph{Evaluation}}

In this chapter we evaluate various aspects of this project, discussing what we
have learned, what went well and what not.
We outline our results, the status of product's requirements and its limitations.
%What we have learned, what went well and what that could have been done better.
Additionally, we talk about social dynamics within the group and with involved parties.

\section{The product}

We were happy with our final product and satisfied with the amount of work we could actually accomplish.
This was thanks to having agreed with the customer on the scope of the problem and by having re-used a
lot of open-source code. We were definitely happy with the project's task which led us to study trending topics
like eHealth and self-quantification. Although we had no previous experiences in designing and developing a web API,
we found it interesting and enjoyable. This resulted in a proactive participation by the group
in the improvement and proposal of requirements, a process which in turn made the project even more interesting
for us to work on.

\subsection{Results}

This project involved designing and developing a number of interoperable systems consisting of:
\begin{itemize}
\item an integration platform
\item a web front-end
\item two Android applications
\item a web application using JSP
\end{itemize}
to analyze the feasibility for a national scale eHealth project to leverage the wide availability of devices
like smartphones and other proprietary solutions like HealthVault.

We have designed and implemented an integration platform as a Spring based web application.
This application exposes a JSON based API through which is possible to exchange data models
representing health measurements such as heart rate and weight.
The web application is deployed using Apache Tomcat and implements persistence using a MySQL database.

%\clearpage
The remaining systems have been developed to demonstrate:
\begin{enumerate}[a.]
\item the functionality of the integration platform by using its API to store and retrieve health measurements
\item the feasibility of integrating with already existing third-party solutions such as Microsoft's HealthVault
\item the difficulty for third-party developers to use the API to make their applications interoperable with the integration platform
\end{enumerate}

Our results outline not only the possibility to leverage handheld devices, which we have demonstrated
using the \textit{Heart rate} application, but also the ease with which third-party developers can make
their application interoperable with our implementation of the integration platform.
In fact, this has been accomplished in both Android applications with less than 100 LOC.
Additionally, our \textit{Weight} application and \textit{HealthVault integration service} have shown how
it is possible to integrate information from other solution provides as well.

Our architectural design choices make the platform well suited to be scaled up and
complemented with security practices. The integration platform in fact, consists
of a web application and a database which can be deployed separately, reducing the amount of resources
needed on a single machine. The same applies to the front-end as well because it interacts with the
integration platform using its API just like any other application and therefore does not rely on being
deployed on the same machine. Furthermore, because we are not employing a classical SOAP client-server
architecture but a more lean, REST-like one, the implementation can be considered lightweight in term of
resources and easier for application developers to interact with.

\subsection{Requirements}

We managed to achieve all requirements and collaborated actively with the customer in the proposal
of new features. However, our product is still only a demonstrative prototype and there is
much room for improvement. Table \ref{table:fulfill-req} presents a complete overview the status of requirements.

\begin{table}[h]
\begin{center}
\begin{tabular}{ | c | c | c | c | c | }
  \hline
  \textbf{ID} & \textbf{Category} &\textbf{Complexity} & \textbf{Priority} & \textbf{Status}\\
  \hline\noalign{\smallskip}\noalign{\smallskip}\hline
  FIP1	& NIPEN				& Med	& High & Completed \\
  FIP2	& NIPEN				& Med	& High & Completed \\
  FIP3	& NIPEN				& Med	& High & Completed \\
  FPI4	& NIPEN				& Med	& High & Completed \\
  FW1   & Front-end			& Med	& High & Completed \\
  FW2   & Front-end			& Low	& Low  & Completed \\
  FHR1	& Heart rate app	& High	& High & Completed \\
  FHR2	& Heart rate app	& Low	& High & Completed \\
  FHR3	& Heart rate app	& Med	& High & Completed \\
  FHV1	& Weighter app		& Med	& High & Completed \\
  FHV2	& Weighter app		& Low	& Med  & Completed \\
  FHV3	& Weighter app		& Low	& High & Completed \\

  FHIS1	& HV integration service	& Med   & High & Completed \\
  FHIS2	& HV integration service	& High  & High & Completed \\
  FHIS3	& HV integration service	& Low	& Med  & Completed \\

  NF1 & Documentation		& High 	& High	& Completed \\
  NF2 & Documentation		& Med 	& High	& Completed \\
  NF3 & Open-source			& Low 	& High	& Completed \\
  NF4 & Interoperability	& Med	& High	& Completed \\
  NF5 & Interoperability	& High	& High	& Completed \\
  NF6 & Accessibility		& Low	& Low	& Completed \\

  \hline
\end{tabular}
\end{center}
\caption{Fulfillment of requirements}
\label{table:fulfill-req}
\end{table}


\subsection{Product limitations}
\label{subsec:limitations}

As every software product, our product has some limitations which should be discussed.
What we present here is a list of the most remarkable limitations of our product.
%that we consider should also be the focus for further development.

\textbf{API design}\newline
At the moment the API only supports two types of data models: weight and heart rate measurements.
The design of an API's data models will have a key role in its success and adoption by
third-parties. What we proposed are two very simple models which include minimal information.
Furthermore, the API is not implementing the basic features of persistent storage (CRUD)
which are necessary for this kind of solution. Only create and read operations are supported,
whereas update and delete are totally missing.

\textbf{Security}\newline
No part of the system implements any form of security.
This was done to reduce the complexity of the system due to time constraints and size of the group.
The customer agreed that security was not a requirement of the system, however in a real-world scenario
appropriate security practices must be adopted.
Especially in projects of this type and scale, security will be a critical requirement
to take into consideration when designing the system and the API as well.

\textbf{Multi-user support}\newline
Some parts of the system such as the API data models, internal data structures and even
database tables do have a basic support for more than one user. However, because no authentication
mechanism is implemented, all actions are effectively carried out by a single user.
Naturally, in a real-world scenario that would not be the case and multi-user support will
be necessary.

\textbf{Deployment}\newline
At the moment both NIPEN, the database and the front-end are deployed together on single server machine.
It would be desirable to deploy these on different machines as that would reduce
the amount of resources needed for a single machine especially considering a large userbase.



\section{Development process}
All in all, we were satisfied with our choice of Scrum and how we adapted
it to our specific needs (see \ref{subsec:devprocess}).
Two weeks sprints proved to be just the right duration to achieve
a high-level goal. Additionally, because there were a total of six sprints,
we had enough chances to adjust our plan to accommodate for changes in requirements.
The overhead introduced by Scrum was negligible and this has proven valuable
given also our resource constraints. There were however some minor issues with our process.

\textbf{Time tracking}\newline
Sometimes, although we had tools to conveniently track the the time spent on the project,
not everybody would do that in due time. This has hindered the effectiveness of burndown charts
to provide meaningful information on the project's progress and could have led to planning issues.
It would have helped to get used to report our times at the end of each day.

\textbf{Effort estimates}\newline
It should be noted that we didn't use any empirical method to estimate the amount of effort
required to complete an activity: the scrum master would simply make a guess.
Especially in the beginning when most of the work was field study and the project
was in the early stages, it was difficult to produce meaningful estimates.
However, using an online tool for scrum we managed to keep an eye on previous
estimates and adjust them based on the amount of effort actually required to complete the tasks.
Furthermore as we became familiar with involved technologies and our own productivity,
the estimates became more and more accurate.

\textbf{Planning}\newline
We think that while our planning was good, at some moments our plan maybe have been not.
Especially in the beginning our plan was not very well structured and this was reflected
by a somewhat overcrowded sprint backlog and less effective means to get an overview
of real the progress of the project. However, our planning became increasingly
methodical, thanks also to the fact that we were using an iterative approach which
left a lot of room for improvement.%, and this resulted in?

\section{Social dynamics}
This section discusses social dynamics during the project, both internally and with
involved parties.

\subsection{Group dynamics}
\label{sec:group}

Group dynamics were definitely a critical aspect of this project.
It would be safe to say that due to the nature of the project itself, its success was based primarily on group dynamics.
We were, in fact, a group of people working to deliver a product to another group of people.
It is easy to convince ourselves that people management in a software process is an essential task for a successful product.
Any successful software project relies on successful group dynamics because no project can be
accomplished by a single man.% We understand that people management will be a critical aspect of our process.
%This encompasses a set of practices, precautions and knowledge which we have taken into consideration and tried our
%best to make our own.

\textbf{Motivate}\newline
Team's commitment will greatly improve many variables in a software project.
People are inevitably influenced by other people around them, especially when they work together.
If any member has a low commitment to the project his attitude will influence negatively other team members.

%A good project manager will motivate the team to work well and gratify members who do so.
%It is a duty of the project manager to instill a positive and enthusiastic attitude in the team.

%To achieve this task, it is important to identify personality types among a group so that
%motivation can be more effective. These can be divided into three categories:
%  \begin{itemize}
%  \item Self-oriented: the work has the purpose to achieve personal interestes.
%  \item Task-oriented: the motivation is related to the task itself.
%  \item Interaction-oriented: the motivation is related to social interaction.
%  \end{itemize}

%It is desiderable that within a group such personality types are balanced.

In our case, the scrum master tried motivating the group by rewarding good jobs
and outlining the importance of the project itself as a chance to acquire knowledge,
to get a good grade (academic achievement) and to do something that stimulated further
research and work in the field.

Feedback from the supervisor and the customer has been constructive and
contributed to further motivate the group.

\textbf{Cooperation}\newline
People should be encouraged to work together, share ideas and thoughts.
We should avoid having one team member whose role is predominant on others'.

During meetings, both internal and external, we had one team member
act as a meeting leader, guiding the customer, supervisor or other team members through
the meeting agenda. However, the meeting leader made sure to involve all team members
in the conversation, asking their opinion on the current topic being discussed.
When the supervisor or the customer asked questions, he made sure not to be the only one
answering by inviting other team members to do so.

\textbf{Size matters}\newline
The fact that the group consisted of three members definitely played a factor in both group dynamics
and project management. Though many may see it as a weakness we believe in our case it was not,
or rather it proved more a strength than a weakness.
In software development, just like in any other engineering practice assigning many people
at a task is unlikely to produce better results than a sufficiently small but motivated group.
%Just like it takes one month to ferment wine, no matter how many people are involved.
Although having a small team means that flexibility in group composition is absent and
the technical skill pool is reduced it is important to realize that a small team
can be more easily motivated and achieve a higher degree of cooperation.
To do so, care should be taken into not letting any member of the group become
a 'black sheep'. Team members should respect each others and their respective work,
nobody should feel less important or cut out of the decision making process.

\textbf{Distance matters}\newline
During the project one team member moved to Oslo and got a full-time job.
This important event required some adjustment to our project management strategy,
especially regarding planning. We thought that the easiest way for him to continue
to contribute to the project was to concentrate on documentation.
This has proved to be generally a good idea although it would have helped
to be more specific about what parts of the documentation he should have focused on.


\subsection{Customer relations}

%\textbf{Relations}\newline
We were definitely pleased with our relations with the customer.
Since our first meeting he was able to explain what would have been the purpose of this project and why
it was relevant for him. He was able to capture our interest and left the task 'open' enough to be stimulating for us.
We collaborated in thinking about features that could have made the product more valuable for him.
The customer was a nice person, smiling and available. He gratified us when he thought we were doing a good job and
that has motivated us to keep it up and work hard. We were glad to be able to work with him.

\textbf{Communication}\newline
Since our customer was located in Oslo, we didn't have many chances to meet in person,
in fact that happened only once. Nevertheless we did have meetings on Skype on a regular weekly basis
with some exceptions due to his work schedule or personal issues.
This wasn't a problem and we were able to manage our work pretty well and autonomously
even those weeks we didn't hear from him.

\subsection{Supervisor relations}

We are happy with our relations with the supervisor.
He made his role in the project clear to us from the beginning so that we knew
what we could count on him for and what not.
We reported him the status of our project weekly and he would always have some feedback
on our progress and also make sure to have a proper understanding of our plans.
His feedback was precious to our project and certainly helped us not to miss any important details.

\textbf{Communication}\newline
Our communication with the supervisor was regular, on a weekly basis.
We had meetings in person on Mondays, agendas being sent one day earlier by mail.
Probably that didn't leave much time to the supervisor to review our documentation.
Additionally, in the beginning we didn't understand very well what kind of documents
we needed to provide him.
However, we feel that these have hardly been problems as we always gave him a
good amount of information on our process and the project's progress and we eventually
become methodical in producing the documentation he needed.

 


%------------------------------
%	THESIS CONTENT - APPENDICES
%------------------------------
\addtocontents{toc}{\vspace{2em}}
\appendix

% Appendix A - Customer Assignment

\chapter{Assignment for Customer driven project} 

\label{AppendixA}
\lhead{Appendix A. \emph{Assignemnt}}

\textbf{Assignment for Customer driven project}

\textbf{Title:} 	National Integration Platform for Citizen Centric eHealth in Norway \\*
\textbf{Customer (Company):} 	The Directorate of Health, Department of the Health Portal\\*
\textbf{Address:}			Universitetsgata 2, Oslo \\


\textbf{Assignment text:} \\
\textit{Background} \\
The Directorate of Health has a national the task of digitalizing Norwegian healthcare, both by providing coordinated services for specialist healthcare (hospitals) and by providing digital services for citizens in general and patients specifically. 
Examples of such services are ePrescriptions, that is implemented on a national basis, the National Summare Care Record, that will go live in Trondheim in August 2013, and the citizen centric health portal  (helsenorge.no) that has been live since June 2011.

National eHealth projects are complex, long running and costly. There are obvious reasons for this. 
Among these are the complexity and criticality of healthcare, and the scale that national eHealth services represents. 

At the same time, the trends in technology development and consumer adaption of new technology continue to develop. 
Moderate prices and consumer friendly devices that monitor individuals’ health and wellness are increasingly becoming available in the market space. 
Combined with a continuous increase in digital competence in the population, they will influence citizens’ behavior and perspective on their own health situation in the future.

In addition to this, private providers develop great eHealth solutions with consumer and patient orientation. Medhelp.org and Healthvault.com are only two among many examples. 
Ambient assisted living has the potential of revolutionizing life for senior citizens with failing health.

The relevant question is: How can the substantial and long running eHealth projects of the government sector connect to and leverage the dynamics in the market and consumer behavior? 
The answer under investigation is the National Integration Platform (NIP) for Citizen Centric eHealth in Norway. 

\textit{Assignment} \\* 
The assignment is to plan, design and describe a NIP, and to develop a prototype.

The task such a platform should fulfil is to offer interoperability with third party solutions based on available application programmable interfaces (APIs).
All third party solution providers must adhere to specified and standardized rules regarding authentication, security model, messaging and privacy to interact with the NIP. 

The intention of such a platform is to enable the following:
\begin{itemize}
\item Citizens’ ability to publish information they produce from devices in their possession and third party software solutions, including smart phone and tablet apps, into the government run citizen centric health portal (helsenorge.no)
\item Citizens’ ability to fetch information about themselves from helsenorge.no to import it into third party software solutions of their own choosing
\end{itemize}

The assignment is to describe the architecture and major components of the NIP, how it will function on the “outside” regarding third party integration, and on the “inside” regarding integration with helsenorge.no.

It is also essential that the solution adhere to Norwegian privacy regulation and information security. 
Its requirements for integration should also encourage privacy by design within third party solutions.

The prototype should make use of one or more use cases to demonstrate how interaction is performed, how privacy and security concerns are managed and how the end user experience will be. \\


Contact details: 

Name:		Helge T. Blindheim\\
Mobile:		466 75 321\\
E-mail:		helgetb@helsedir.no

 % A Assignment 
\chapter{How to build the project}
\label{AppendixB}
\lhead{Appendix B. \emph{Build project}} 


Tomcat 

Android app:
The APK can be found here: https://github.com/Andersos/NIPEN/tree/master/nipen-heart-rate % B How to build
\chapter{Templates}
\label{AppendixC} 
\lhead{Appendix C. \emph{Templates}}

In this appendix we presents our templates for our meeting documents
and status reports.

\section{Weekly status report}

\begin{center}
Weekly status report \#X\\*
Week NN \\*
Dates 2013-MM-DD - 2013-MM-DD \\*
TDT4290 Customer Driven Project - Group 17
\end{center}

\textbf{1.2 Meetings}

\textbf{2. Plan for next week}

\textbf{3. Milestones}

\textbf{4. Problems}

\textbf{1. Work done}

\iffalse
\begin{table}[H]
\begin{center}
\begin{tabular}{ l | l | l }
  \hline
  Activity & Planned & Actual \\
  \hline\noalign{\smallskip}\noalign{\smallskip}\hline
  Studies & Number & Number \\
  Project management & Number & Number \\
  System developement & Number & Number \\
  Application development & Number & Number \\
  Database developement & Number & Number \\
  Testing & Number & Number \\
  Report & Number & Number \\
  \hline
\end{tabular}
\end{center}
\caption{Activity chart}
\label{table:activityChartStatusReport}
\end{table}
\fi

\clearpage
\section{Advisor meeting agenda}

\begin{center}
Advisor meeting agenda \#X\\
2013-MM-DD\\
\end{center}

\begin{table}[H]
\begin{center}
\begin{tabular}{ l | l }
Time and Data & 2013-MM-DD HH:MM \\
Place & Room \\
Attendees & Full name \\
Referent & Full name \\
\end{tabular}
\end{center}
%\caption{Activity chart}
%\label{table:activityChartAdvisorAgenda}
\end{table}

1 Work done \\
1.1 Meetings \\
2 Plan for next week \\
3 Milestones \\
4 Problems \\

\clearpage
\section{Advisor meeting notes}

\begin{center}
Advisor meeting notes \#X \\
2013-MM-DD \\
\end{center}

\begin{table}[H]
\begin{center}
\begin{tabular}{ l | l }
Time and Data & 2013-MM-DD HH:MM \\
Place & Room \\
Attendees & Full name \\
Referent & Full name \\
\end{tabular}
\end{center}
%\caption{Activity chart}
%\label{table:activityChartAdvisorAgenda}
\end{table}

Topics discussed
\begin{itemize}
\item topic
\item topic
\end{itemize}
 
\clearpage
\section{Customer meeting notes}

\begin{center}
Customer meeting notes \#X \\
2013-MM-DD \\
\end{center}

\begin{table}[H]
\begin{center}
\begin{tabular}{ l | l }
Time and Data & 2013-MM-DD HH:MM \\
Attendees & Full name \\
Referent & Full name \\
\end{tabular}
\end{center}
%\caption{Activity chart}
%\label{table:activityChartAdvisorAgenda}
\end{table}

Topics discussed
\begin{itemize}
\item topic
\item topic
\end{itemize}

\clearpage
\section{Internal meeting notes}

\begin{center}
Internal meeting notes \#X \\
2013-MM-DD \\
\end{center}

\begin{table}[H]
\begin{center}
\begin{tabular}{ l | l }
Time and Data & 2013-MM-DD HH:MM \\
Attendees & Full name \\
Referent & Full name \\
\end{tabular}
\end{center}
%\caption{Activity chart}
%\label{table:activityChartAdvisorAgenda}
\end{table}

Topics discussed
\begin{itemize}
\item topic
\item topic
\end{itemize} % C Templates
\chapter{Meeting documents}
\label{AppendixD}
\lhead{Appendix D. \emph{Meeting documents}}

In this appendix we present some documents produced for supervisor meetings.

\iffalse
\section{Overview of advisor meetings}

\begin{table}[H]
\begin{center}
\begin{tabular}{ l | l | l }
  \hline
  Weekday & Date & Meeting Number \\
  \hline\noalign{\smallskip}\noalign{\smallskip}\hline
  Wednesday 	&	21.08.2013 & 1 \\
  	% Was there not a advisor meeting on the first friday or something?
	Monday		&	26.08.2013 & 2 \\
	Monday		& 02.09.2013 & 3 \\
	Monday 		&	09.09.2013 & 4 \\
	Monday 		&	16.09.2013 & 5 \\
	Monday 		&	23.09.2013 & 6 \\
	Monday 		&	30.09.2013 & 7 \\
	Monday 		&	07.10.2013 & 8 \\
	Monday 		&	14.10.2013 & 9 \\
	Monday 		&	21.10.2013 & 10 \\
	Monday 		&	28.10.2013 & 11 \\
	Monday 		&	04.11.2013 & 12 \\
	Monday 		&	11.11.2013 & 13 \\
	Monday 		&	18.11.2013 & 14 \\
  \hline
\end{tabular}
\end{center}
\caption{Customer meetings}
\label{table:customermeetings}
\end{table}
\fi


\section{Weekly status reports}
In this section we present some of the status reports we produced.

\clearpage
% WEEK 35
\subsection{Week 35}

\begin{center}
Weekly status report \#01\\*
Week 35 \\*
Dates 2013-08-26 - 2013-09-01 \\*
TDT4290 Customer Driven Project - Group 17
\end{center}

\textbf{1. Work done}

This week was mostly focused on initial planning of the project.
This included but was not limited to:
\begin{itemize}
\item Internal meeting planning
\item Customer meeting planning
\item Iteration planning
\item Milestone planning
\item Preliminary studies
\end{itemize}


The team agreed the following weekly meetings:
\begin{itemize}
\item On Mondays from 10 to 19
\item On Wednesdays from 15 to 19
\end{itemize}

Each team member agreed to work another 8 hours per week by himself.
The team worked on a preliminary system architecture that was submitted
to the Customer in order to make clear to him our understanding of the project goal.
The Customer was then contacted on Monday the 26th through Skype and provided positive feedback.
The Customer and the team agreed on when to have the next meetings.

A Google Drive folder was created in order to share and collaborate in creating and editing documents. 
The Customer has also access to this folder and has already used it to have a look at the preliminary system architecture.

The team also planned iterations (their number, duration and focus) and milestones for the project
and initiated preliminary studies on relevant technologies.

The team chose a LaTex template for the report.

\textbf{1.1 Meetings}

Monday
from 10 to 19, all members present.
Skype call with the customer, meeting with the supervisor.

Wednesday
from 15 to 19, all members present.

\textbf{2. Plan for next week}

Based on the feedback received from our Customer, we have identified the scope of the project and will conduct studies on relevant technologies in the project’s domain. We will try to identify some formal requirements
for the system, or rather parts of it, based on some use cases and proceed to plan the work for the next iteration.

\textbf{3. Milestones}

Milestones reached: none.

\textbf{4. Problems}

The team is still unsure about the commitment of its fourth member
who still has to show up. Furthermore the timetable for the course
is not clear and inconsistent between sources.

\newpage
% WEEK 36
\subsection{Week 36}

\begin{center}
Weekly status report \#02\\*
Week 36 \\*
Dates 2013-09-02 - 2013-09-08 \\*
TDT4290 Customer Driven Project - Group 17
\end{center}

\textbf{1. Work done}

The team studied relevant existing solutions and technologies.
In particular, among the ones we took into consideration there is an integration platform suggested by the customer himself.

The team set-up the working environments and other tools such as:
\begin{itemize}
\item LaTeX environments for writing the report
\item a server machine for product deployment
\item a git repository managed with GitHub
\item Apache Tomcat web server
\item Android’s SDK
\item IDEs
\end{itemize}

During the Skype meeting, the customer mentioned Microsoft’s  HealthVault and expressed his wish to
see it supported by the integration platform (the product).

The team developed a prototype Android application to monitor the user’s heart rate. An assessment
of the risks for this project was performed and will be updated weekly or once per iteration.

\textbf{1.1 Meetings}

Monday
from 10 to 19, all members present.
Skype call with the customer, meeting with the supervisor.

Wednesday
from 15 to 19, all members present. Emanuele joined two hours late because he had to attend a lecture. 
He had notified the other team members beforehand.

\textbf{2. Plan for next week}

Because team members are not familiar with the technologies involved  additional study is required in order
to decide upon system architecture.
Regarding HealthVault, the team still needs to asses the feasibility of supporting such technology in the product. 
For the next iteration the team will start working on an initial system based on the Spring Framework, the
Apache Tomcat server and database backend.  The team will also start working on the report. 
We will produce more use cases and define requirements.

\textbf{3. Milestones}

Milestones reached: none

\textbf{4. Problems}

The customer is not providing us some documents we requested.
These are not critical for the moment.

\newpage
% WEEK 37
\subsection{Week 37}

\begin{center}
Weekly status report \#03\\*
Week 37 \\*
Dates 2013-09-09 - 2013-09-15 \\*
TDT4290 Customer Driven Project - Group 17
\end{center}

\textbf{1. Work done}

During this week the team has worked toward an initial implementation of the system. 
The database server has been installed and configured on the server machine. The team also started
working on the API itself and on the portal web-page. 
A initial set of requirements for part of the system has been identified. Risk analysis has been updated.

\textbf{1.2 Meetings}

Monday
from 10 to 19, all members present.
Skype call with the customer, meeting with the supervisor.

Wednesday
from 15 to 19, all members present.

\textbf{2. Plan for next week}

Further studies need to be conducted on relevant technologies. 
The team will meet the customer in person on Friday. 
We plan to demonstrate some basic system functionality by then. 
The team should try to get started on the mid-term report as soon as possible.

\textbf{3. Milestones}

Milestones reached: none

\textbf{4. Problems}

Given the size of our group, we expect writing the report will delay our progress with the product itself.





\iffalse

\newpage
% WEEK 38
\subsection{Week 38}

\begin{center}
Weekly status report \#04\\*
Week 38 \\*
Dates 2013-09-16 - 2013-09-22 \\*
TDT4290 Customer Driven Project - Group 17
\end{center}

\textbf{1. Work done}

The team has finished developing a first prototype of the system which has been shown to the customer. 
He provided very positive feedback.
In order to help us plan and prioritize our work we asked the Customer about what he’d like us to start working on next. 
The Customer expressed his interest for HealthVault's integration. 
We informed the Customer that during the next iteration we will mostly focus on report work anyway.

Risk analysis has been updated. Milestones have been reviewed based on supervisor’s advice.

\begin{table}[H]
\begin{center}
\begin{tabular}{ l | l | l }
  \hline
  Activity & Planned & Actual \\
  \hline\noalign{\smallskip}\noalign{\smallskip}\hline
  Studies & 25 & 5 \\
  Project management & 12 & 12 \\
  System developement & 25 & 50 \\
  Application development & 10 & 6 \\
  Database developement & 10 & 2 \\
  Testing & 6 & 4 \\
  Report & 20 & 0 \\
  \hline
\end{tabular}
\end{center}
\caption{Activity chart}
\label{table:activityChartStatusReportWeek38}
\end{table}

\textbf{1.2 Meetings}

Monday from 10 to 19
Emanuele and Sebastian present.
Skype call with the customer, meeting with the supervisor.

Wednesday from 15 to 19
Emanuele and Sebastian present.

Friday from 12.30 to 14.30
Emanuele and Sebastian present, Anders joined on Skype.
Meeting with the Customer in person. Demonstration of the product.

\textbf{2. Plan for next week}

During the next week the team will focus primarily on the report.
We shared our plan for the next iteration with the Customer, as we do regularly. 
Emanuele will also study HealthVault in order to assess the possibility of its integration in our platform. 
Because we are confident we can use HealthVault  we have planned some work for its integration as well.

\textbf{3. Milestones}

M1 First prototype of the system.

\begin{enumerate}
\item An Android application which measures the user’s heart rate and forwards it to the integration platform.
\item A Spring REST endpoint that handles such data and stores it in a local MySQL database.
\item A webpage frontend displaying the data using charts.
\end{enumerate}

Associated deliverables: none


\textbf{4. Problems}

None.


\newpage
% WEEK 39
\subsection{Week 39}

\begin{center}
Weekly status report \#05\\*
Week NN \\*
Dates 2013-09-23 - 2013-09-29 \\*
TDT4290 Customer Driven Project - Group 17
\end{center}

\textbf{1. Work done}

During last week the group has worked exclusively on the report.

\textbf{1.2 Meetings}

Monday from 10 to 19
Emanuele and Sebastian present. Anders joined on Skype.
Meeting with the supervisor. Unable to contact the customer.

Wednesday from 15 to 19
Emanuele and Sebastian present. Anders joined on Skype.


\textbf{2. Plan for next week}

During next week the group will still focus on report work but also spend a little bit of time on development.

\textbf{3. Milestones}

Milestones achieved: None

\textbf{4. Problems}

None.

\fi





\clearpage
\section{Advisor meeting agenda}

See below some of the agendas for our meetings with the advisor.

\subsection{Week 35}

\begin{center}
Agenda for advisor meeting \#01\\*
2013-08-26\\*
\end{center}

\begin{table}[H]
\begin{center}
\begin{tabular}{ l | l }
Time and Data & 2013-08-26 12:00 \\
Place & BU1-108 \\
\end{tabular}
\end{center}
\end{table}

\begin{enumerate}[1.]
\item Approval of agenda
\item Approval of minutes from last advisor meeting. See attached meetings notes.
\item Comments to last weeks minutes
\item Approval of the weekly report. See attached weekly report.
\item Topics to be discussed
\item Progress done last week and plan for the next iteration
\item Other
\item For next meeting
\end{enumerate}

\newpage
\subsection{Week 36}

\begin{center}
Agenda for advisor meeting \#02\\*
2013-09-02\\*
\end{center}

\begin{table}[H]
\begin{center}
\begin{tabular}{ l | l }
Time and Data & 2013-09-02 12:00 \\
Place & BU1-108 \\
\end{tabular}
\end{center}
\end{table}


\begin{enumerate}[1.]
\item Approval of agenda
\item Approval of minutes from last advisor meeting. See attached meetings notes
\item Comments to last weeks minutes
\item Approval of the weekly report. See attached weekly report.
\item Topics to be discussed
\item Progress done last week and plan for the next iteration.
\item Other
\item For next meeting
\end{enumerate}

\clearpage
\section{Advisor meeting notes}

See below some of the notes taken during our meetings with the advisor.

\subsection{Week 45}

\begin{center}
Advisor meeting notes \#X \\
2013-11-04 \\
\end{center}

\begin{table}[H]
\begin{center}
\begin{tabular}{ l | l }
Time and Data & 2013-11-04 14:00 \\
Place & IT building \\
Attendees & Emanuele Di Santo, Sebastian Zalewski \\
Referent & Sebastian Zalewski \\
\end{tabular}
\end{center}
%\caption{Activity chart}
%\label{table:activityChartAdvisorAgenda}
\end{table}

Topics discussed
\begin{enumerate}[1.]
\item The supervisor asks about the meeting with the customer last week. We tell him that the meeting was over skype and not in person. We tell him that the customer was happy. The supervisor asks if the customer asked about a user manual to the application. We tell him that we are going to add it to the report.

\item We ask if it is possible to only let one person present. It can work if only one person talks in the presentation, but the supervisor thinks that all should contribute. It implies that everybody has contributed, if everyone talks. The supervisor tells that it can work if one person talks about the slides, while the other person shows the demonstration. The presentation usually consists of 30 min slides 15 minutes for questions (usually 5-10 minutes questions). Common questions: How we distributed the work, which parts to improve etc… The presentation is the best way to present the work for the censor. The censor might not look thoroughly through the report.

\item The supervisor asks about the status of the report. We tell him that the report is around 100 pages (the main part), the supervisor thinks that is enough
\end{enumerate}

\clearpage
\subsection{Week 46}
\begin{center}
Advisor meeting notes \#X \\
2013-11-11 \\
\end{center}

\begin{table}[H]
\begin{center}
\begin{tabular}{ l | l }
Time and Data & 2013-11-11 14:00 \\
Place & IT building \\
Attendees & Emanuele Di Santo, Sebastian Zalewski \\
Referent & Sebastian Zalewski \\
\end{tabular}
\end{center}
\end{table}

The supervisor has read throught the report.
He thinks that it is good. Content is interesting, language is good.

\textbf{1. Feedback on the report}\newline
-Section 1.2 only have for example name of the customer in title.\\
-Title of section 1.7 “reoprt structure” instead of “summary”.\\
-Move the project tools to the last section of the chapter. (Maybe move to the prestudies).\\
Project management should only contain final tools used (does not need an explanation why).
- Weekly meetings is not so interesting in project manager. It is however possible to keep it
where it is. (Combine with quality assurance, (change name of quality assurance for example management rules))\\
- Functional requirements: give an overview of the whole system. Use natural language to give an
overview how the whole application works together. One paragraph, 3 or 4 sentences how everything works together.\\
- Add a use case diagram for the integration platform as well. Integration platform is not a good name
(in requirement), it doesn’t look like an individual system.\\
- We need to explain why we choose the architecture in the arhictecture chapter.
Constraints. Not highest priority, just to make perfect if we have time. Should contains drivers.\\
- The supervisor wants to add the security chapter into another chapter for example evaluation or prestudies.\\
- Should mention the burndown charts in the text.\\
- Tell something about technical issues, changes in requirements etc. in sprint chapters. A seperate section for it.\\
- Sprint backlog should be moved to the appendix, because it might not be so interesting for the censor.\\
- Too many things in the conclusion. Should cover mainly two parts: the evaluation of the product, and an
evaluation of the process and management.\\
- We should have an evaluation chapter as well as conclusion.\\
- The evaluation: an evaluation how we did it.\\
- Conclusion: summary of the whole report, future work.\\
- The reflection can be simpel (1 or 2 pages) improvement of the course\\
- Should make the headings more descriptive\\

\textbf{2. Questions}\newline
We ask the supervisor what he thinks about the tables. He thinks that all the tables are in
good quality with usefull information. The gant diagram: font might be a bit small (scale it better).
The class diagram should contain alll the classes instead of detailed information of methods and attributes.
Not necceassry to add more to the appendices.
Should explain what JSON is in footnote or something. Not neccessary to add to preliminary studies.

\textbf{3. Delivery}\newline
The customer wants a copy of the report on the 17th for final comments. % D Advisor 
\chapter{Customer meeting documents}
\label{AppendixE}
\lhead{Appendix E. \emph{Customer}}

\section{Agenda for customer meeting}

\section{Minutes of customer meeting}
 % E Customer
\chapter{Internal meeting notes}
\label{AppendixF}
\lhead{Appendix F. \emph{Internal}}


\section{Overview of internal meetings}

\begin{table}[h]
\begin{center}
\begin{tabular}{ l | l | l }
  \hline
  Weekday & Date & Meeting Number \\
  \hline\noalign{\smallskip}\noalign{\smallskip}\hline
  	Wednesday 	&	21.08.2013 & 1 \\
	Monday		&	26.08.2013 & 2 \\
	Monday		& 	02.09.2013 & 3 \\
	Monday 		&	09.09.2013 & 4 \\
	Monday 		&	16.09.2013 & 5 \\
	Monday 		&	23.09.2013 & 6 \\
	Friday 		&	29.09.2013 & 7 \\
	Monday 		&	30.09.2013 & 8 \\
	Monday 		&	07.10.2013 & 9 \\
	Monday 		&	14.10.2013 & 10 \\
	Monday 		&	21.10.2013 & 11 \\
	Monday 		&	28.10.2013 & 12 \\
	Monday 		&	04.11.2013 & 13 \\
	Monday 		&	11.11.2013 & 14 \\
	Monday 		&	18.11.2013 & 15 \\
  \hline
\end{tabular}
\end{center}
\caption{Internal meetings}
\label{table:internalmeetings}
\end{table}

\section{Agenda for internal meeting}

\section{Minutes of internal meeting} % F Internal
\chapter{Test results}
\label{AppendixG}
\lhead{Appendix G. \emph{Test results}}

\section{Functional tests}

\begin{table}
\begin{center}
\begin{tabular}{ | l | p{10cm} | }
	\hline
	\textbf{Test}	&	\textbf{ID 1} \\
	\hline\noalign{\smallskip}\noalign{\smallskip}\hline
	Application name	& NIPEN \\
	Name				& Heart rate REST (read) \\
	Requirement			& \hyperref[table:reqip]{FIP1} \\
	Description			& Test the REST endpoint for receiving heart rate data models \\
	Preconditions		& 	\par The IP is deployed and running on a server machine 
							\par The test is run on the server machine or alternatively
							one that has access to the server and whose address is replaced to the
							string \verb|localhost| in the test.
							\par The machine on which the test is run has the \verb|curl| program installed.
							\par At least one heart rate measurement has been stored on the IP \\
	Steps 				&	1. Run the following:
							\begin{verbatim}
							curl localhost:8080/nipen/human/api/heart_rates
							\end{verbatim}
							\\
	Postconditions		& A JSON valid, comma separated list of heart rate data models consistent with 
							entries in the databased hosted on the server. \\
	Results				& An array of heart rate JSON strings is returned. 
						  The values are consistent with the values on the database. 
						  Hence, the test was successful. \\
	Comments			& - \\
	Status				& OK \\
	Tester				& Sebastian Zalewski \\
	Date				& 08-11-2013 \\
	\hline
\end{tabular}
\end{center}
\end{table}

\begin{table}
\begin{center}
\begin{tabular}{ | l | p{10cm} | }
	\hline
	\textbf{Test}	&	\textbf{ID 2} \\
	\hline\noalign{\smallskip}\noalign{\smallskip}\hline
	Application name	& NIPEN \\
	Name				& Weight REST (read) \\
	Requirement			& \hyperref[table:reqip]{FIP2} \\
	Description			& Test the REST endpoint for receiving weight data models \\
	Preconditions		&	\par The IP is deployed and running on a server machine
							\par The test is run on the server machine or alternatively
							one that has access to the server and whose address is replaced to the
							string \verb|localhost| in the test.
							\par The machine on which the test is run has the \verb|curl| program installed.
							\par At least one weight measurement has been stored on the IP \\
	Steps 				&	1. Run the following:
							\begin{verbatim}
							curl localhost:8080/nipen/human/api/weights
							\end{verbatim}
							\\
	Postconditions		& A JSON valid, comma separated list of heart rate data models consistent with 
							entries in the database hosted on the server. \\
	Results				& An array of weight JSON strings is returned. 
						  The values are consistent with the values on the database. 
						  Hence, the test was successful. \\
	Comments			& - \\
	Status				& OK \\
	Tester				& Sebastian Zalewski \\
	Date				& 08-11-2013 \\
	\hline
\end{tabular}
\end{center}
\end{table}

\begin{table}
\begin{center}
\begin{tabular}{ | l | p{10cm} | }
	\hline
	\textbf{Test}	&	\textbf{ID 3} \\
	\hline\noalign{\smallskip}\noalign{\smallskip}\hline
	Application name	& NIPEN \\
	Name				& Heart rate REST (write) \\
	Requirement			& \hyperref[table:reqip]{FIP3} \\
	Description			& Test the REST endpoint for requesting heart rate data models \\
	Preconditions		&	\par The IP is deployed and running on a server machine
							\par The test is run on the server machine or alternatively
							one that has access to the server and whose address is replaced to the
							string \verb|localhost| in the test.
							\par The machine on which the test is run has the \verb|curl| program installed. \\
	Steps 				&	1. Run the following \begin{verbatim}
							curl -X POST -H "Content-Type: application/json" 
							-d "{\"userId\":1,\"timestamp\":
							\"2013-11-08 17:33:05\",\"value\":60,
							\"unit\":\"bpm\"}" 
							localhost:8080/nipen/api/human/heart_rate
							\end{verbatim} \\
	Postconditions		& A database entry coherent with the JSON data submitted is created on the server. \\
	Results				& A new entry on the database is created with the sent values.
						  Hence, the test was successful. \\
	Comments			& - \\
	Status				& OK \\
	Tester				& Sebastian Zalewski \\
	Date				& 08-11-2013 \\
	\hline
\end{tabular}
\end{center}
\end{table}

\begin{table}
\begin{center}
\begin{tabular}{ | l | p{10cm} | }
	\hline
	\textbf{Test}	&	\textbf{ID 4} \\
	\hline\noalign{\smallskip}\noalign{\smallskip}\hline
	Application name	& NIPEN \\
	Name				& Weight REST (write) \\
	Requirement			& \hyperref[table:reqip]{FIP4} \\
	Description			& Test the REST endpoint for requeting weight data models \\
	Preconditions		&	\par The IP is deployed and running on a server machine
							\par The test is run on the server machine or alternatively
							one that has access to the server and whose address is replaced to the
							string \verb|localhost| in the test.
							\par The machine on which the test is run has the \verb|curl| program installed. \\
	Steps 				&	1. Run the following \begin{verbatim}
							curl -X POST -H "Content-Type: application/json" 
							-d "{\"userId\":1,\"timestamp\":
							\"2013-11-08 17:45:05\",\"value\":89,
							\"unit\":\"kg\"}" 
							localhost:8080/nipen/api/human/weight
							\end{verbatim} \\
	Postconditions		& A database entry coherent with the JSON data submitted is created on the server. \\
	Results				& A new entry on the database is created with the sent values.
						  Hence, the test was successful. \\
	Comments			& - \\
	Status				& OK \\
	Tester				& Sebastian Zalewski \\
	Date				& 08-11-2013 \\
	\hline
\end{tabular}
\end{center}
\end{table}

\begin{table}
\begin{center}
\begin{tabular}{ | l | p{10cm} | }
	\hline
	\textbf{Test}	&	\textbf{ID 5} \\
	\hline\noalign{\smallskip}\noalign{\smallskip}\hline
	Application name	& Heart Rate Application \\
	Name				& Heart rate measurement \\
	Requirement			& \hyperref[table:reqheartrate]{FHR1} \\
	Description			& Test the heart rate measurement functionality \\
	Preconditions		& The application has started \\
	Steps 				&	\par 1. User holds his finger on the camera applying a slight pressure \\
	Postconditions		&	\par 1. The icon on the left side of the screen is blinking 
							\par 2. A measurement appears after 4 seconds at most \\
	Results				& The result was a heart measurement of 77, which is accurate.
						  However, the measurement didn't appear until after 6 or 7 seconds.
						  This might be because of wrong finger placement, and may also vary from phone to phone.
						  The measurement was quite accurate and hence the test was successful. \\
	Comments			&	The measurement is expected to be rough.
							Any value between 60 and 100 is okay as long as it actually varies slightly based
							on the tester's perceived heart rate.  \\
	Status				& OK \\
	Tester				& Sebastian Zalewski \\
	Date				& 08-11-2013 \\
	\hline
\end{tabular}
\end{center}
\end{table}

\begin{table}
\begin{center}
\begin{tabular}{ | l | p{10cm} | }
	\hline
	\textbf{Test}	&	\textbf{ID 6} \\
	\hline\noalign{\smallskip}\noalign{\smallskip}\hline
	Application name	& Heart Rate Application \\
	Name				& Send heart rate \\
	Requirement			& \hyperref[table:reqheartrate]{FHR3} \\
	Description			& Test the heart rate application interoperability \\
	Preconditions		&	\par 1. The application has started
							\par 2. A heart rate measurement has been acquired \\
	Steps 				&	\par 1. User presses \textbf{Send} button \\
	Postconditions		&	\par 1. The application shows a toast
							\par 2. The measurement is acquired by the Integration Platform \\
	Results				& The measurement was sent to the front-end and contained the right values.
						  The test was successful. \\
	Comments			& - \\
	Status				& OK \\
	Tester				& Sebastian Zalewski \\
	Date				& 08-11-2013 \\
	\hline
\end{tabular}
\end{center}
\end{table}


\begin{table}
\begin{center}
\begin{tabular}{ | l | p{10cm} | }
	\hline
	\textbf{Test}	&	\textbf{ID 7} \\
	\hline\noalign{\smallskip}\noalign{\smallskip}\hline
	Application name	& Weigh Application \\
	Name				& Data fetching from HealthVault\\
	Requirement			& \hyperref[table:reqweight]{FHV1} \\
	Description			& Test HealthVault connectivity \\
	Preconditions		&	\par 1. The application has started
							\par 2. The user has authenticated to HealthVault \\
	Steps 				&	\par 1. User clicks on the \textit{Fetch} button \\
	Postconditions		&	\par 1. The data is correctly fetched from HealthVault
							\par 2. A list is populated with the data \\
	Results				& The values appearing on the app are corresponding to the values on HealthVault.
						  Thus the test was successful. \\
	Comments			& - \\
	Status				& OK \\
	Tester				& Sebastian Zalewski \\
	Date				& 08-11-2013 \\
	\hline
\end{tabular}
\end{center}
\end{table}

\begin{table}
\begin{center}
\begin{tabular}{ | l | p{10cm} | }
	\hline
	\textbf{Test}	&	\textbf{ID 8} \\
	\hline\noalign{\smallskip}\noalign{\smallskip}\hline
	Application name	& Weight Application \\
	Name				& Sending data to HealthVault through app \\
	Requirement			& \hyperref[table:reqweight]{FHV3} \\
	Description			& Test the HealthVault application data forwarding functionality \\
	Preconditions		&	\par 1. The application has started
							\par 2. The application has fetched some data from HealthVault \\
	Steps 				&	\par 1. User enters a weight value
							\par 2. User clicks on the \textit{Send} button \\
	Postconditions		&	\par 1. A success message is showed
							\par 2. Data is correctly acquired by the integration platform \\
	Results				& The data sent is acquired by the integration platform, with the right values.
						  This is seen on the front-end and in the database.
						  The test was successful. \\
	Comments			& - \\
	Status				& OK \\
	Tester				& Sebastian Zalewski \\
	Date				& 08-11-2013 \\
	\hline
\end{tabular}
\end{center}
\end{table}

\begin{table}
\begin{center}
\begin{tabular}{ | l | p{10cm} | }
	\hline
	\textbf{Test}	&	\textbf{ID 9} \\
	\hline\noalign{\smallskip}\noalign{\smallskip}\hline
	Application name	& HealthVault Integration Service \\
	Name				& Sending data to HealthVault through web service \\
	Requirement			& \hyperref[table:reqwebservice]{FHIS3} \\
	Description			& Test if the web service i capable of sending values to HealthVault \\
	Preconditions		&	\par 1. The user has logged in to HealthVault through the web service \\
	Steps 				&	\par 1. Enter a value into the weight field
							\par 2. Click on submit to HealthVault button \\
	Postconditions		&	\par 1. The submitted value is stored on HealthVault \\
	Results				& The value entered is registered on HealthVault.
						  The test was successful. \\
	Comments			& - \\
	Status				& OK \\
	Tester				& Sebastian Zalewski \\
	Date				& 08-11-2013 \\
	\hline
\end{tabular}
\end{center}
\end{table}

\begin{table}
\begin{center}
\begin{tabular}{ | l | p{10cm} | }
	\hline
	\textbf{Test}	&	\textbf{ID 10} \\
	\hline\noalign{\smallskip}\noalign{\smallskip}\hline
	Application name	& HealthVault Integration Service \\
	Name				& Polling service \\
	Requirement			& \hyperref[table:reqwebservice]{FHIS1} and \hyperref[table:reqwebservice]{FHIS2} \\
	Description			& Test if new weight measurements stored on HealthVault are sent to NIPEN \\
	Preconditions		&	\par 1. The user has logged in to HealthVault through the web service
							\par 2. The polling service is disabled \\
	Steps 				&	\par 1. User clicks on \textit{Enable} button, to start the polling 
							\par 2. User adds a value into HealthVault \\
	Postconditions		&	\par 1. The value is sent to the integration platform
							\par 2. The sent value can be viewed on the front-end \\
	Results				& The value that was entered on HealthVault is visible on the front-end.
						  Hence, the test was successful. \\
	Comments			& - \\
	Status				& OK \\
	Tester				& Sebastian Zalewski \\
	Date				& 08-11-2013 \\
	\hline
\end{tabular}
\end{center}
\end{table} % Test results

\addtocontents{toc}{\vspace{2em}}
\backmatter

%-----------------
%	BIBLIOGRAPHY
%-----------------
\label{Bibliography}
\lhead{\emph{Bibliography}} % Change the page header to say "Bibliography"
\bibliographystyle{unsrtnat} % Use the "unsrtnat" BibTeX style for formatting the Bibliography
\bibliography{Bibliography} % The references (bibliography) information are stored in the file named "Bibliography.bib"
\end{document}  